% Modernised reading — fonds Alexandre Grothendieck, shelfmark 146, pages 1-133.
%
% This is an INTERPRETATION, not a transcription. It re-reads the pages in
% today's notation and vocabulary, states the hypotheses the manuscript leaves
% implicit, and drops the critical apparatus so that the mathematics reads
% straight through. Everything the brackets and underlines carried moves into a
% footnote. It is held to being mathematically correct as it stands; where the
% manuscript is loose or mistaken the true statement appears and the footnote
% says what the page has. In any dispute of substance the transcription
% governs: whoever cites Grothendieck cites batch-01.fr.tex ... batch-07.fr.tex.
%
% Scope: seven batches (1-20, 21-40, 41-60, 61-80, 81-100, 101-120, 121-133),
% all transcribed, read whole before any of this was written, and written as
% ONE file for the whole shelfmark. Many pages are blank versos or unrelated
% administrative paper; the transcriptions' headers list them.
%
% The spine. The folder builds discrete models S_0T_G of the "Teichmüller
% groupoids" S T_G = Pi_1(SP_G) — fundamental groupoids of the torsors of unit
% tangent directions (at the marked points and at the nodes) over the strata
% M_G of the compactified moduli spaces — by inducing them on finite sets of
% real "base points at infinity", and relates the strata by generalisation
% morphisms S T_G -> S T_{G'}. Stations:
%   - pp. 1-14: S_0T_G for an MD graph of modular dimension 0, via the
%     groupoids E(eps) attached to an extension and a torsor;
%   - pp. 15-17: card I = 4, the tangent bundles over M^_{0,4} = P^1_J;
%   - pp. 18-54: the line P(V_C) marked by a regular polygon, its tautological
%     C*-bundle minus n lines (the fair draft pp. 42-54 is read FIRST, then
%     the computations of pp. 19-39 that use it);
%   - pp. 55-77: working sheets (icosahedral census for (0,5), base points at
%     infinity, Belyi maps of the monogon and bigon, M^_{0,5}, line bundles
%     L_i, platonic groups);
%   - pp. 78-98: "Groupoïdes" — presentations of extensions and of groupoids
%     over groupoids, pi-fibrations, (pi, mu)-fibrations;
%   - pp. 99-124: S_0T_G in general, generalisation morphisms by holomorphic
%     surgery, modular dimensions 0 and 1;
%   - pp. 125-132: "Graphes MD", the combinatorics used from page 2 on.
%
% Notation fixed for the whole folder. Arcs (half-edges) hat-vec-A, involution
% sigma, origin map "or" (pages 127-130 write sigma for the origin and
% sigma-bar for the involution); free edges I, bound edges A, all edges hat-A.
% Teichmüller groupoids written \mathcal{T} (pages 102-124 trace a Gothic T).
% U is the circle group; base-point sets are eps (pages 89-96 write E, pages
% 122-123 a calligraphic E); extensions are \mathcal{E}; the 2-groups acting
% on base points are gamma. TWO base-point sets share the name S_0 on the
% pages: eps(G) (one orientation per vertex, pp. 10-14, 2^{|S|} points) and
% S eps_G (one real direction per edge, pp. 111-116, 2^{|hat A|} points); the
% groupoid induced on the second is written S_0T'_G here.
%
% Substantive departures from the pages, each footnoted where it occurs:
%   - p. 13: prod_{Delta in diag Q} Delta is the set of SIDES of the square,
%     not of its vertices;
%   - pp. 16-17: \hat T_i = O(1) is right in EGA's convention V(E); the
%     tangent line has degree -1 (psi_i has degree 1);
%   - p. 35: Lemma 1 is false without its (illegible) orientation hypothesis;
%     counterexample given;
%   - p. 37: sin(pi t), not sin t; p. 39: Lemma 4 needs Delta a vector line,
%     and the right-hand vertical arrow is L_{-x, -phi(x)};
%   - p. 64: the page's normalisations are checked, g_1 and g_2 are Belyi maps;
%   - p. 66: a scribbled ratio is not reproduced (its sign does not check);
%   - p. 80: the proof that one-sided conjugation relations suffice is ours;
%   - p. 86: the system defines T~, not T;
%   - p. 88: the last group of the exact sequence is Aut_C; fullness needed;
%   - p. 94: (17) is an equivalence for POINTED groupoids (or on iso classes);
%   - p. 109: ]0,1]^C, as the page's own addition says;
%   - p. 114: S eps_G has 2^{card hat-A} elements, not 2^{card A};
%   - p. 116: relation (2) follows from (1);
%   - p. 121: a square structure is a (2,2) partition, not (1,2);
%   - p. 122: T_{J_II} = U^{J_II/sigma};
%   - p. 127: (6) nu = card I; (9) is the dimension of the TYPE, the modular
%     dimension of the graph is 3g-3+nu-c(G) (what (10) and pp. 111, 117 use).
%
% Announced and never established in the folder: the expression of
% S_0T_{0,I}, card I = 4, by generators and relations (pp. 26-32 eliminate
% generators and stop); the end of the argument "Alors on sait qu'on a" (p.
% 62); the inverse construction of p. 97; the proof that (12)-(14) define T~
% (p. 86); the transition condition between fibres (p. 95); the filtered
% colimit of p. 99; the explicit relations for S_0T^Q_G (p. 123, "j'y
% reviendrai"); the equianharmonic case (p. 124 stops on "On").
%
% Pass: Opus 5.5 (claude-opus-5-5), 2026-10-10 — first pass, unchecked against the pages by a human.
\documentclass[11pt,a4paper]{article}
\input{../preamble/grothendieck}

\folder{146}
\pages{1}{133}
\dating{[à partir de 1978]}
\watermark{Édition de démonstration\\interprétation personnelle de l'œuvre}
\foldertitle{Groupoïdes de Teichmüller et les S0 Tg,!ν, S0 TG : notes manuscrites (s.d.) — lecture modernisée du dossier entier}

\begin{document}

\section*{Résumé}

\begin{resume}
Une surface — une sphère, un tore, une surface à plusieurs anses — peut être
déformée de mille façons, et l'ensemble de toutes ses formes possibles est
lui-même un espace, l'\emph{espace des modules}. Un chemin fermé dans cet
espace ramène la surface à sa forme de départ, mais en l'ayant tordue en
route : c'est ainsi qu'apparaissent les symétries profondes d'une surface,
celles que forment les \emph{groupes modulaires} (ou groupes de Teichmüller).
Ces cent trente-trois pages cherchent une façon de décrire ces groupes, et les
« groupoïdes » qui les rassemblent, par des données finies et explicites.

L'idée qui les gouverne est de regarder les surfaces au bord de leur espace
des modules, là où elles dégénèrent. Une surface très dégénérée se casse en
morceaux élémentaires, des sphères à trois trous — des « pantalons », dit-on
aujourd'hui —, recollés deux à deux ; le plan de ce recollement est un graphe,
que le dossier appelle \emph{graphe MD}. Un pantalon n'a presque aucune forme
propre : trois points sur une sphère sont toujours dans la même position.
Il ne lui reste qu'une donnée, une petite boussole en chacun de ses trois
points, la direction tangente. Le dossier fait de ces boussoles le cœur du
calcul : en chaque point, deux directions réelles privilégiées, et pour aller
de l'une à l'autre un demi-tour de l'aiguille. Faire tourner l'aiguille d'un
tour complet, c'est exactement le « twist » qui tord la surface autour de ce
point. Les pages 1 à 14 décrivent ainsi, par une algèbre d'entiers pairs et
impairs, le groupoïde attaché à une surface entièrement découpée en pantalons.

Viennent ensuite les surfaces un peu moins dégénérées, qui ont un seul degré
de liberté : la sphère à quatre points, dont l'espace des formes est lui-même
une sphère privée de trois points. Les pages 15 à 54 l'étudient avec un soin
presque artisanal — un hexagone régulier, ses symétries, les deux pôles où la
configuration est la plus symétrique — pour choisir des points de base que
toutes les symétries respectent, et écrire les chemins entre eux par
générateurs et relations. Les pages 99 à 124 expliquent enfin comment passer
d'une surface dégénérée à ses voisines moins dégénérées : on enlève de petits
disques autour des nœuds et l'on recolle les bords en les tournant d'un angle
choisi — une « chirurgie holomorphe ». Entre-temps, un cahier intitulé
« Groupoïdes » met au point l'outil général (pages 78 à 98) : comment
présenter un groupe ou un groupoïde construit par extension, et comment un
petit ensemble de points de base, stable par un groupe de symétries, suffit à
reconstituer tout le groupoïde par une recette purement algébrique.

Des feuillets de travail s'intercalent : un recensement des configurations de
cinq points sur la sphère, qui fait surgir l'icosaèdre et ses 12 sommets, 30
arêtes et 20 faces ; les équations du « monogone » et du « bigône »,
$-\frac14\bigl(z - \frac1z\bigr)^2$, saluées d'un « de merveille ! » ; une
première formulation des « points base à l'infini » sur un corps quelconque.
La datation vient des versos : un sujet d'examen de juin 1978, des imprimés de
1979, une circulaire du 12 février 1982 ; le dossier s'étend donc au moins
jusqu'en 1982.

Où cela mène-t-il ? C'est le chantier de la « tour de Teichmüller » que
l'\emph{Esquisse d'un programme} (1984) présentera comme un jeu de Lego : tout
s'engendrerait à partir des surfaces de dimension modulaire $0$ et $1$, avec
des relations venues de la dimension $2$ — la page 111 désigne déjà ces cas
comme « les plus intéressants ». La notion de point base tangentiel a été
formalisée par Deligne en 1989 ; que la tour des groupes modulaires
s'engendre ainsi à partir des dimensions $1$ et $2$ a été démontré par
Hatcher, Lochak et Schneps en 2000.

\keywords{Teichmüller groupoid, mapping class group, tangential base point, Deligne–Mumford compactification, stable graph, boundary stratum, plumbing construction, Dehn twist, action groupoid, central extension of groupoids, group presentation, pure braid group, hyperplane arrangement, Belyi map, dessin d'enfant, cartographic group, icosahedral group, equivariant Picard group, Lego–Teichmüller game, pants decomposition}
\end{resume}

\pagerange{1}{133}
\section*{Le fil du dossier, et les conventions}

\subsection*{Les stations}

\begin{itemize}
\item \textbf{Calcul des $S_0\mathcal{T}_G$ en dimension modulaire nulle
(pages 1 à 14)}, sous une chemise titrée de sa main (page 1) : les groupoïdes
$\mathcal{E}(\varepsilon)$ d'une extension et d'un torseur ; le cas d'un
sommet, $S_0\mathcal{T}_{0,I}$ avec $\operatorname{card} I = 3$ ; le groupe
fondamental « variant » ; le cas d'un graphe ; l'exemple du carré.
\item \textbf{Le cas $\operatorname{card} I = 4$ (pages 15 à 17)} : les fibrés
tangents aux points marqués au-dessus de $\widehat{M}_{0,I} \simeq
\mathbb{P}^1_J$.
\item \textbf{La droite projective marquée par un polygone régulier (pages 18
à 54)}, sous une chemise (page 18) : « Étude de $P\Sigma_J^{*}$ (et de
$P\Sigma_n^{*}$\ldots) ». Une rédaction au propre, numérotée de 1 à 7 de sa
main (pages 42 à 54), fixe le cadre ; les calculs des pages 19 à 39
(générateurs et relations, lemmes d'homotopie) s'en servent. On lit donc la
rédaction d'abord.
\item \textbf{Feuillets de travail (pages 55 à 77)} : configurations de cinq
points et icosaèdre ; points base à l'infini ; monogone et bigône ;
$\widehat{M}_{0,5}$ ; fibrés en droites $L_i$ ; groupes des polyèdres réguliers.
\item \textbf{« Groupoïdes » (pages 78 à 98)}, sous une chemise (page 78) :
présentations d'extensions de groupes et de groupoïdes ; $\pi$-fibrations ;
$(\pi, \mu)$-fibrations et groupoïdes induits sur des points de base.
\item \textbf{$S_0\mathcal{T}_G$ en général, et les morphismes de
généralisation (pages 99 à 124)} : réductions ; chirurgie holomorphe ;
dimension modulaire $0$, puis $1$.
\item \textbf{« Graphes MD » (pages 125 à 132)}, sous une chemise (page 125) :
le tableau des graphes de petits types, les définitions, les opérations.
\end{itemize}

Les définitions combinatoires dont tout le dossier se sert ne viennent qu'à la
fin (pages 127 à 132) ; on les rappelle ci-dessous, et la dernière section les
reprend dans l'ordre des pages.

\subsection*{Les conventions}

\emph{Graphes MD.} Un graphe est donné par un ensemble fini $S$ de sommets, un
ensemble fini $\hat{\vec A}$ d'\emph{arcs} (demi-arêtes), une application
origine $\mathrm{or} : \hat{\vec A} \to S$, une involution $\sigma$ de
$\hat{\vec A}$, et un genre $g : S \to \mathbb{N}$\footnote{Les pages 127 à 130
notent $\sigma$ l'application origine et $\bar\sigma$ l'involution ; les pages
1 à 14 et 99 à 124 notent $\sigma$ l'involution et « or » l'origine (page 14).
On suit ces dernières. Les pages 1 à 14 écrivent $\vec{\hat A}$ pour
$\hat{\vec A}$.}. Les points fixes de $\sigma$ sont les arêtes \emph{libres},
$I = \hat{\vec A}^{\sigma}$ — les points marqués ; les autres arcs forment
$\vec A$, et $A = \vec A/\sigma$ est l'ensemble des arêtes \emph{liées} — les
nœuds ; $\hat A = \hat{\vec A}/\sigma = I \amalg A$. On note $\hat{\vec A}_s$
les arcs d'origine $s$ et $\hat\nu_s = \operatorname{card} \hat{\vec A}_s$
l'ordre du sommet. Le graphe est \emph{MD} (stable au sens de Deligne et
Mumford) si $2g_s + \hat\nu_s \geq 3$ pour tout $s$. Une courbe stable de type
$G$ a pour normalisée la réunion des courbes lisses $\widetilde{X}_s$ de genre
$g_s$ marquées par $\hat{\vec A}_s$, recollées selon $\sigma$.

\emph{Espaces de modules.} $M_{g,I}$ est l'espace (le champ) des modules des
courbes lisses de genre $g$ marquées par $I$, $\widehat{M}_{g,I}$ sa
compactification de Deligne--Mumford\footnote{Les pages écrivent $\hat M$ pour
la compactification, la notation de l'époque, $\overline{M}$ aujourd'hui.},
$M_G$ la strate des courbes de type $G$. La \emph{dimension modulaire} de $G$
est $\dim M_G = \sum_s (3g_s - 3 + \hat\nu_s)$ ; elle est nulle exactement
quand tous les sommets sont de genre $0$ et d'ordre $3$, et alors $M_G$ est un
point.

\emph{Torseurs tangents.} $\mathbb{U} = \{ z \in \mathbb{C} \mid |z| = 1\}$.
Pour une courbe $X$ et un point lisse $x$, $T_x$ est le
$\mathbb{U}$-torseur des directions tangentes en $x$, $(T_xX \smallsetminus
0)/\mathbb{R}^{*}_{+}$ ; pour un nœud $a$, de branches $a', a''$, $T_a =
T_{a'} \wedge T_{a''}$ (produit contracté). Au-dessus d'une strate, ils
forment des fibrés principaux : $SP_G = \prod_{a \in \hat A} T_a$, de groupe
$\mathbb{U}^{\hat A}$ ; $P_G = \prod_{a \in A} T_a$ ; et plus généralement
$P_{G,K} = \prod_{a \in K} T_a$ pour $K \subset \hat A$ (page 103). Les
groupoïdes fondamentaux correspondants sont
\[
S\mathcal{T}_G = \Pi_1(SP_G), \qquad \mathcal{T}_G = \Pi_1(P_G), \qquad
R\mathcal{T}_G = \Pi_1(M_G),
\]
le dernier au sens des orbifoldes ; $S\mathcal{T}_G$ est une extension de
$R\mathcal{T}_G$ par $\mathbb{Z}^{\hat A}$, les générateurs du noyau étant
les tours complets des directions tangentes — en termes d'aujourd'hui, les
twists de Dehn autour des points marqués et des nœuds\footnote{Le titre du
dossier note $\mathcal{T}^{!}_{g,\nu}$ ce que la page 102 identifie à
$S\mathcal{T}_{g,I}$ : pour une courbe lisse, c'est le groupoïde dont les
groupes sont les groupes modulaires de la surface de genre $g$ à $\nu$
composantes de bord, extension du groupe modulaire pur par
$\mathbb{Z}^{\nu}$. Les pages 102 à 124 tracent la lettre $\mathcal{T}$ en
gothique ; ce sont les mêmes objets.}.

\emph{Points de base.} Les indices $0$ à gauche ($S_0\mathcal{T}$,
${}_0\mathcal{T}$) désignent toujours un sous-groupoïde \emph{plein}, induit
sur un ensemble fini de points de base réels, noté $\varepsilon$ ; le
$2$-groupe qui agit simplement transitivement sur lui est noté $\gamma$.
Deux choix de points de base portent le même nom sur les pages, et on les
distingue : l'ensemble $\varepsilon(G)$ des pages 10 à 12, une orientation par
sommet ($2^{\operatorname{card} S}$ éléments), et l'ensemble $S\varepsilon_G$
des pages 113 à 116, une direction réelle par arête ($2^{\operatorname{card}
\hat A}$ éléments). On garde $S_0\mathcal{T}_G$ pour le groupoïde induit sur le
premier, et l'on écrit $S_0\mathcal{T}'_G$ pour celui induit sur le second.

\emph{Extensions.} $\mathcal{E}$ désigne toujours une extension de groupes, et
$\mathcal{E}_z$ sa fibre au-dessus d'un élément $z$ du quotient. L'ensemble à
deux éléments des orientations d'un ensemble fini $I$ est $\omega(I)$ ; $\mu_2
= \{\pm 1\}$.

\pagerange{1}{14}
\section*{Calcul des $S_0\mathcal{T}_G$ en dimension modulaire nulle (pages 1
à 14)}

\pagerange{2}{5}
\subsection*{Le groupoïde d'une extension et d'un torseur (pages 2 à 5)}

Soit $G$ un graphe MD de dimension modulaire nulle. Coupons-le en tous ses
nœuds : on obtient un graphe $G^{b}$ dont chaque sommet $s$, isolé, porte trois
arcs libres, et pour lequel $S_0\mathcal{T}_{G^b} = \prod_{s \in S}
S_0\mathcal{T}_{0, \hat{\vec A}_s}$. Le groupoïde de $G$ s'en déduit par le
changement de groupe d'opérateurs
\[
(1)\qquad S_0\mathcal{T}_G = S_0\mathcal{T}_{G^b} \wedge^{\mathbb{Z}^{\hat{\vec
A}}} \mathbb{Z}^{\hat A},
\]
le long de l'homomorphisme $\mathbb{Z}^{\hat{\vec A}} \to \mathbb{Z}^{\hat A}$
qui somme les coordonnées des deux arcs d'une même arête — la traduction, au
niveau des groupes fondamentaux, de $T_a = T_{a'} \wedge T_{a''}$\footnote{La
page dit « changement de groupes interne », associé à $\vec{\hat A} \to \hat
A$. Une note marginale, au-dessus, n'est lisible que par fragments.}. Tout se
ramène donc au cas d'un seul sommet : une droite projective $\mathbb{P}^1_I$
marquée par un ensemble $I$ à trois éléments.

\emph{Les points de base.} Trois points de $\mathbb{P}^1(\mathbb{C})$ sont sur
un unique cercle ; une orientation $\alpha \in \omega(I)$ de $I$ (un ordre
cyclique) oriente ce cercle, et donne en chaque $i \in I$ un vecteur tangent
unitaire, d'où un isomorphisme de $\mathbb{U}$-torseurs $\varphi_{\alpha,i} :
T_i \simeq \mathbb{U}$, et un point $\bar\alpha$ du $\mathbb{U}^I$-torseur
$P\mathbb{U}_{0,I} = \prod_{i \in I} T_i$. D'où une injection
\[
(3)\qquad \omega(I) \hookrightarrow P\mathbb{U}_{0,I},
\]
compatible avec l'homomorphisme diagonal $\{\pm 1\} \to \{\pm1\}^I \subset
\mathbb{U}^I$ : renverser l'orientation retourne chacun des trois vecteurs.
Tout est compatible avec l'action de $\mathfrak{S}_I$, qui agit sur
$\omega(I)$ par la signature. On \emph{définit} $S_0\mathcal{T}_{0,I}$ comme le
sous-groupoïde plein du groupoïde fondamental de $P\mathbb{U}_{0,I}$ induit sur
$\omega(I)$.

\emph{Le groupoïde d'une extension.} Pour un $\mathbb{U}$-torseur $P$ et $x, y
\in P$, les classes de chemins de $x$ à $y$ sont la fibre
$\widetilde{\mathbb{U}}_{y-x}$ du revêtement universel $\mathbb{R} \to
\mathbb{U}$, la composition étant l'addition. La formule a un sens purement
algébrique. Soit
\[
(8)\qquad 0 \to \pi \to \mathcal{E} \to \Gamma \to 0
\]
une extension de groupes abéliens et $\varepsilon$ un $\Gamma$-torseur. Le
groupoïde $\mathcal{E}(\varepsilon)$ a pour objets $\varepsilon$, pour flèches
\[
(9)\qquad \mathrm{Hom}(x, y) = \mathcal{E}_{y-x},
\]
et se compose par l'addition dans $\mathcal{E}$ ; chaque
$\mathrm{Hom}(x,y)$ est un $\pi$-torseur. C'est le \emph{groupoïde d'action}
de $\mathcal{E}$ agissant sur $\varepsilon$ à travers $\Gamma$ : son ensemble de
flèches est $\mathcal{E} \times \varepsilon$, la source est la seconde
projection, le but l'action\footnote{Le nom n'est pas sur la page, qui écrit
$\mathrm{Fl}\,\mathcal{E}(\varepsilon) = \mathcal{E} \times \varepsilon$ et la
loi $(u, vx) \circ (v, x) = (uv, x)$.}. Si $\varepsilon' \to \varepsilon$ est un
morphisme de torseurs au-dessus d'un homomorphisme $\Gamma' \to \Gamma$, le
groupoïde induit sur $\varepsilon'$ est $\mathcal{E}'(\varepsilon')$, où
$\mathcal{E}' = \mathcal{E} \times_{\Gamma} \Gamma'$ est l'extension image
inverse (11) : en effet $\mathrm{Hom}(x', y') = \mathcal{E}_{f(y')-f(x')}$, et
$f(y') - f(x')$ est l'image de $y' - x'$.

Appliqué à l'injection (3), cela donne un isomorphisme de
$\mathbb{Z}^I$-groupoïdes
\[
(12)\qquad S_0\mathcal{T}_{0,I} \simeq \mathcal{E}_I\bigl(\omega(I)\bigr),
\]
où $\mathcal{E}_I$ est l'image inverse de l'extension $0 \to \mathbb{Z}^I \to
\mathbb{R}^I \to \mathbb{U}^I \to 0$ par l'homomorphisme diagonal $\mu_2 \to
\mathbb{U}^I$.

\pagerange{6}{10}
\subsection*{L'extension $\mathcal{E}_I$ et le groupe fondamental variant
(pages 6 à 10)}

Explicitement, $\mathcal{E}_I$ est formé des $v \in (\frac12\mathbb{Z})^I$ dont
les coordonnées sont toutes entières ou toutes demi-entières. Mesurons les
chemins en \emph{demi-tours} : en multipliant par $2$,
\[
\mathcal{E}_I \simeq \{\, n \in \mathbb{Z}^I \mid \text{les } n_i \text{ ont
tous la même parité} \,\},
\]
le noyau $\mathbb{Z}^I$ (les tours complets) s'y plongeant comme $2\mathbb{Z}^I$
et la projection sur $\mu_2$ étant la parité commune. C'est l'extension
obtenue à partir de
\[
(14)\qquad 0 \to \mathbb{Z} \xrightarrow{\;2\;} \mathbb{Z} \to
\mathbb{Z}/2\mathbb{Z} \simeq \mu_2 \to 0
\]
par l'homomorphisme diagonal $\mathbb{Z} \to \mathbb{Z}^I$ (diagramme (15)).
Ainsi, pour $x, y \in \omega(I)$,
\[
(16)\qquad \mathrm{Hom}_{S_0\mathcal{T}_{0,I}}(x, y) =
\begin{cases}
\{\, n \in \mathbb{Z}^I : \text{tous les } n_i \text{ pairs} \,\} & \text{si } x = y,\\
\{\, n \in \mathbb{Z}^I : \text{tous les } n_i \text{ impairs} \,\} & \text{si } x \neq y,
\end{cases}
\]
la composition étant l'addition : c'est la « description algébrique
complète » de $S_0\mathcal{T}_{0,I}$ en termes du seul ensemble $I$\footnote{La
page écrit $\mathrm{Hom}(x,y) = T_{x-y} \wedge^{\mathbb{Z}} \mathbb{Z}^I$, où
$T_{x-y}$ est le $\mathbb{Z}$-torseur des entiers pairs ou impairs, et
l'identifie à $\mathbb{Z}^I$ ou à $\mathrm{Imp}^I$ (entiers impairs) ; le
$\mathbb{Z}^I$ y agit par $n \mapsto n + 2m$. C'est la même chose.}.

\emph{Le groupe fondamental variant.} $\mathfrak{S}_I$ agit sur
$S_0\mathcal{T}_{0,I}$, et l'on forme, pour $x \in \omega(I)$,
\[
(17)\qquad \pi_1(S_0\mathcal{T}_{0,I}, \mathfrak{S}_I; x) =
\{\, (\alpha, l) \mid \alpha \in \mathfrak{S}_I,\ l : \alpha x \to x \,\},
\]
avec la loi
\[
(18)\qquad (\alpha, l)(\alpha', l') = \bigl(\alpha\alpha',\ l \circ
\alpha(l')\bigr).
\]
C'est exactement la construction $\pi_1(X, \Gamma, a)$ que le dossier 141
(page 14) donne pour un groupe agissant sur un espace, le groupe fondamental
orbifold d'aujourd'hui, ici pour l'action de $\mathfrak{S}_I$ sur le
groupoïde\footnote{Le rapprochement est de cette édition ; les deux lois de
composition coïncident.}. Un automorphisme $\alpha$ de $\mathbb{P}^1_I$
envoie $T_i$ sur $T_{\alpha(i)}$ et le vecteur de l'orientation $x$ sur celui
de $\alpha x$, en respectant le sens de rotation ; dans les coordonnées (16),
il agit donc sur $\mathbb{Z}^I$ en permutant les coordonnées. D'où
\[
\pi_1(S_0\mathcal{T}_{0,I}, \mathfrak{S}_I; x) \simeq G(I) = \bigl\{\,
(\alpha, n) \in \mathfrak{S}_I \times \mathbb{Z}^I \bigm| n_i \equiv
\tfrac{1 - \mathrm{sgn}(\alpha)}{2} \bmod 2 \text{ pour tout } i \,\bigr\},
\]
avec $(\alpha, n)(\alpha', n') = (\alpha\alpha', n + \alpha\cdot n')$.

\textbf{Proposition (pages 8 à 10).} $G(I)$ est l'extension de
$\mathfrak{S}_I$ par $\mathbb{Z}^I$ obtenue en prenant l'image inverse de (14)
par la signature $\mathfrak{S}_I \to \mu_2$, puis l'image directe par
l'homomorphisme diagonal $\mathbb{Z} \to \mathbb{Z}^I$ ($\mathfrak{S}_I$
agissant sur $\mathbb{Z}^I$ par permutation). Elle n'est pas scindée ; sa
restriction à $\mathfrak{S}_I^{+}$, le stabilisateur de $x$, est canoniquement
scindée par $\alpha \mapsto (\alpha, 0)$.

En effet, l'application $\bigl(v, (\alpha, m)\bigr) \mapsto (\alpha, 2v +
m\mathbf{1})$, où $m \equiv \frac{1-\mathrm{sgn}(\alpha)}{2}$, est un
homomorphisme surjectif du produit semi-direct de $\mathbb{Z}^I$ par l'image
inverse de (14) sur $G(I)$, de noyau les $(-k\mathbf{1}, (1, 2k))$ : c'est la
première assertion. Si $\sigma$ est une transposition fixant $k$, et $(\sigma,
n) \in G(I)$, alors $(\sigma, n)^2 = (1, n + \sigma n)$ a pour coordonnée $k$
l'entier $2n_k$, non nul puisque $n_k$ est impair : aucun relèvement de
$\sigma$ n'est d'ordre $2$, et l'extension, déjà restreinte au sous-groupe
d'ordre $2$ engendré par $\sigma$, n'est pas scindée\footnote{La page le
montre en projetant sur la coordonnée $k$, ce qui redonne (14), visiblement non
scindée ; c'est le même argument.}. Enfin, pour $\alpha$ pair, $n = 0$ est
permis.

La page s'étonne : « Il est un peu étrange » que ce groupe ne dépende pas du
choix de $x$. Il se construit sans référence à $x$ ; $\mathfrak{S}_I$ agit
donc sur lui de façon compatible, bien qu'aucun scindage ne décrive cette
action — sur $\mathfrak{S}_I^{+}$ elle est celle du scindage canonique. C'est,
pour la partie $\mathfrak{S}_3$ du groupe $\Gamma_{0,3}$, l'extension
$S\Gamma$ de $\Gamma_{0,3}$ par $\mathbb{Z}^{\{0,1,\infty\}}$ que le dossier
145 (page 31) laisse à expliciter\footnote{Rapprochement de cette édition ;
le dossier 145 adjoint en outre la conjugaison complexe, absente ici.}.

\pagerange{10}{14}
\subsection*{Un graphe de dimension modulaire nulle ; l'exemple du carré
(pages 10 à 14)}

Pour $G$ de dimension modulaire nulle, posons
\[
\varepsilon(G) = \prod_{s \in S} \omega(\hat{\vec A}_s), \qquad
\gamma(G) = \{\pm 1\}^S,
\]
et soit $\mathcal{E}(G^b)$ l'extension de $\gamma(G)$ par $\mathbb{Z}^{\hat{\vec
A}}$ déduite de la puissance $S$-ième de (14) par l'homomorphisme
$\mathbb{Z}^S \to \mathbb{Z}^{\hat{\vec A}}$, $n \mapsto (n_{\mathrm{or}(h)})_h$,
puis $\mathcal{E}(G)$ l'extension de $\gamma(G)$ par $\mathbb{Z}^{\hat A}$
déduite de celle-ci par la somme $\mathbb{Z}^{\hat{\vec A}} \to
\mathbb{Z}^{\hat A}$ (24)--(26). Alors
\[
(27)\qquad S_0\mathcal{T}_G \simeq \mathcal{E}(G)\bigl(\varepsilon(G)\bigr) :
\qquad \mathrm{Ob} = \varepsilon(G), \quad \mathrm{Hom}(x,y) =
\mathcal{E}(G)_{y-x},
\]
et la page propose de prendre (28) pour \emph{définition} de
$S_0\mathcal{T}_G$\footnote{Les numéros (22) à (26) sont récrits et lus avec
peine ; (21) n'apparaît pas.}.

\emph{L'exemple du carré.} Le graphe à deux sommets reliés par une arête $a$,
chacun portant deux arêtes libres, correspond à un ensemble $I$ de quatre
points muni d'une partition de type $(2,2)$, qu'on peut voir comme les deux
diagonales d'un carré combinatoire $Q$ de sommets $I$. Alors $S \simeq
\mathrm{diag}\,Q$, l'arc de la diagonale $\Delta_s$ vers l'arête $a$ étant
noté $\tilde s$ ; $\hat{\vec A}_s = \Delta_s \amalg \{\tilde s\}$, et une
orientation de cet ensemble à trois éléments équivaut au choix du successeur
de $\tilde s$, d'où $\omega(\hat{\vec A}_s) \simeq \Delta_s$ et
\[
\varepsilon(G) \simeq \prod_{\Delta \in \mathrm{diag}\,Q} \Delta,
\]
torseur sous $\gamma(G) = \{\pm1\}^{\mathrm{diag}\,Q}$. Un élément de ce
produit est un couple de sommets pris sur les deux diagonales, c'est-à-dire
deux sommets adjacents : $\varepsilon(G)$ est l'ensemble des quatre
\emph{côtés} de $Q$. L'élément diagonal $(-1,-1)$ fait passer au côté opposé,
et le quotient, $\bigwedge_{\Delta} \Delta$, est l'ensemble à deux éléments des
paires de côtés opposés\footnote{La page écrit « $A(Q)$ l'ens. des quatre
sommets, s'identifiant à $\prod_{\Delta \in \mathrm{diag}\,Q} \Delta$ », puis
« $I = A^{\mathrm{lib}} = A(Q)$ ens. des sommets de $Q$ », et nomme ce quotient
« codiag$(Q)$ ». Le produit des deux diagonales est l'ensemble des côtés, pas
celui des sommets ; la lettre $A$ (arêtes) le suggère d'ailleurs. Identifier
$\varepsilon(G)$ à $I$ demanderait un choix.}. On obtient $S_0\mathcal{T}_G$ à
partir de l'extension $0 \to \mathbb{Z}^S \xrightarrow{2} \mathbb{Z}^S \to
\{\pm1\}^S \to 0$ et du torseur $\varepsilon(G)$, puis par l'homomorphisme
$\mathbb{Z}^S \to \mathbb{Z}^{\hat A} \simeq \mathbb{Z}^I \times \mathbb{Z}$,
dont les deux composantes sont l'homomorphisme $\mathbb{Z}^S \to \mathbb{Z}^I$
déduit par contravariance de l'application $I \to S$ (de degré $2$) et la somme
$\mathbb{Z}^S \to \mathbb{Z}$, le facteur $\mathbb{Z}$ étant celui de l'arête
$a$ (page 14).

\pagerange{15}{17}
\section*{Le cas $\operatorname{card} I = 4$ (pages 15 à 17)}

Soit $I$ à quatre éléments, $J$ l'ensemble à trois éléments de ses partitions
de type $(2,2)$, $\Sigma = \Sigma_I = \mathbb{P}^1_J$ la droite projective
marquée par $J$ et $\omega = \omega(J)$. On a canoniquement
\[
(4)\qquad M_{0,I} \simeq \Sigma^{*} = \Sigma \smallsetminus J, \qquad
\widehat{M}_{0,I} \simeq \Sigma :
\]
les trois points du bord de $\widehat{M}_{0,I}$ sont les trois façons de
grouper les quatre points deux par deux. $\mathfrak{S}_I$ agit à travers
$\mathfrak{S}_I \to \mathfrak{S}_J$, de noyau le groupe de Klein $V(I)$, dont
les trois éléments non triviaux correspondent aux éléments de $J$. Les deux
points de $\Sigma$ fixés par $\mathfrak{S}_J^{+}$ (les configurations
équianharmoniques) sont dans les deux hémisphères que découpe le cercle réel
$\Sigma_{\mathbb{R}}$ passant par $J$ ; d'où un isomorphisme
\[
(2)\qquad \omega \simeq \Sigma^{\mathfrak{S}_J^{+}},
\]
qui associe à un ordre cyclique sur $J$ l'hémisphère dont le bord, orienté par
cet ordre, est $\Sigma_{\mathbb{R}}$\footnote{La page dessine la sphère, son
pôle $P$, l'équateur et un parallèle ; les légendes et une note marginale sont
en partie illisibles.}.

La question des pages 16 et 17 est de savoir comment les fibrés tangents
$T_i$ aux points marqués se prolongent à $\widehat{M}_{0,I}$. La page écrit
$\Sigma_J = \mathbb{P}(V(J))$, où $0 \to V(J) \to \mathcal{O}^J
\xrightarrow{\text{somme}} \mathcal{O} \to 0$, et trouve, après avoir comparé
les deux possibilités en un point où la courbe a une structure de carré,
\[
\widehat{T}_i \simeq \mathcal{O}_\Sigma(1) \qquad \text{pour chaque } i \in I,
\]
sans torsion par $\omega$. Il faut le lire dans la convention de l'époque, où
un fibré vectoriel $\mathbb{V}(\mathcal{F})$ est nommé par le faisceau de ses
formes linéaires : le faisceau conormal de la section $i$ est
$\mathcal{O}(1)$, et le fibré tangent, comme faisceau de sections, est de
degré $-1$. C'est ce qu'on sait aujourd'hui : la classe $\psi_i$ est de degré
$1$ sur $\overline{M}_{0,4}$\footnote{La page écrit $\widehat{T}_i \simeq
\mathcal{O}(1)^{\otimes\,?} \otimes_{\mathbb{Z}} L(\omega)$, l'exposant incertain,
avec l'alternative $L(\omega) = \mathbb{Z}$ ou $\mathbb{Z}(\omega)$, et conclut
$L(\omega) = \mathbb{Z}$ par comparaison en la section $\alpha'$. La lecture
« conormal de degré $1$ » est de cette édition ; elle est confirmée par la
page 51 et la marge de la page 52, où le $\mathbb{C}^{*}$-torseur de
$\mathcal{O}_\Sigma(1)$ est identifié au plan privé de l'origine, mais avec
l'action \emph{opposée} des homothéties.}.

Comme les quatre torseurs sont canoniquement un même torseur, leur produit est
déduit de lui par l'homomorphisme diagonal $\mathbb{G}_m \to
\mathbb{G}_m^I$ :
\[
SP_I \simeq \bigl(\mathbb{V}^{*}(\mathcal{O}(1)) \,\big|\, \Sigma^{*}\bigr)
\wedge^{\mathbb{G}_m} \mathbb{G}_m^I, \qquad
S\mathcal{T}_{0,I} \simeq \Pi_1\bigl(\mathbb{V}^{*}(\mathcal{O}(1)) \,\big|\,
\Sigma^{*}\bigr) \wedge^{\mathbb{Z}} \mathbb{Z}^I,
\]
où $\mathbb{V}^{*}$ désigne le fibré privé de sa section nulle\footnote{La
première formule de la page porte $\check{\mathbb{V}}^{*}(\mathcal{O}(1))^I$,
l'exposant $I$ lu avec doute.}. Tout se ramène donc au groupoïde fondamental
d'un seul $\mathbb{C}^{*}$-torseur au-dessus de la sphère privée de trois
points. Ce torseur est le plan $V^{*}$ privé de trois droites par l'origine
(pages 51 et 52 ci-dessous) ; son groupe fondamental est le groupe de tresses
pures $P_3 \simeq F_2 \times \mathbb{Z}$, le facteur central $\mathbb{Z}$
étant le tour de la fibre, et l'on retrouve que le groupe modulaire pur de la
sphère à quatre composantes de bord est $F_2 \times
\mathbb{Z}^4$\footnote{Ces identifications sont de cette édition. Le plan
$V(J)$ des vecteurs de somme nulle de $\mathbb{C}^J$ est l'espace des
configurations de trois points de $\mathbb{C}$ à translation près, et les
trois droites ôtées sont les diagonales $z_j = z_k$.}.

La chemise de la page 18 annonce la suite : « Étude de $P\Sigma_J^{*}$ (et de
$P\Sigma_n^{*}$\ldots) $S_0\mathcal{T}_{0,4}$, $S_0\mathcal{T}_{0,I}$
($\operatorname{card} I = 4$) ».

\pagerange{41}{54}
\section*{La droite projective marquée par un polygone régulier (pages 41 à
54)}

Cette rédaction, numérotée 1 à 7 de sa main, est le cadre au propre de ce que
les pages 19 à 39 calculent ; on la lit en premier. La page 41 n'est qu'un
décalque pâle de la page 42.

\pagerange{42}{45}
\subsection*{Le polygone et son groupe diédral (pages 42 à 45)}

Soient $n \geq 3$, $m = 2n$, et un polygone régulier à $m$ côtés dans un plan
vectoriel réel $V$, d'ensemble de sommets $\tilde J$\footnote{Une note
marginale dit que $n \geq 2$ suffirait, le cas $n = 2$ étant « très
dégénéré ».}. La donnée de $(V, \tilde J)$ équivaut à celle d'un torseur sous
le groupe diédral $D_m = \mu_m \rtimes \mathbb{Z}/2$, groupe d'automorphismes
du modèle $(\mathbb{C}, \mu_m)$, où $\mu_m$ agit par homothéties et
$\mathbb{Z}/2$ par conjugaison. On pose $D = \mathrm{Aut}(V, \tilde J)$,
$\underline{\omega}$ l'ensemble des deux orientations de $V$, $D^{+}$ le
sous-groupe des rotations, noyau de $D \to \mathrm{Aut}(\underline{\omega})$.
Les deux rotations d'angle $2\pi/m$ dans l'un et l'autre sens forment une
copie de $\underline{\omega}$ dans $D^{+}$ ; ce sont deux générateurs inverses
l'un de l'autre, $\omega' = \omega^{-1}$. Les sommets sont notés
$\widetilde{R}_i$, $i \in \tilde J$.

Comme $m$ est pair, $\tilde J$ est stable par $a_0 : x \mapsto -x$, seul
élément d'ordre $2$ du groupe cyclique $D^{+}$. Pour $j \in \tilde J$, la
réflexion $r_j$ fixant $\widetilde{R}_j$ fixe aussi $\widetilde{R}_{a_0(j)}$ :
il y a $n$ telles réflexions, la moitié de $D \smallsetminus D^{+}$, les $n$
autres étant les réflexions par les milieux des côtés opposés. Il y a sur $V$
une unique forme quadratique $D$-invariante valant $1$ sur les sommets, qui
fait de $V$ un plan euclidien.

Soit $\Sigma = \mathbb{P}(V_{\mathbb{C}})$ la droite projective complexe
associée au complexifié. $D$ y agit à travers $D/\{1, a_0\}$, groupe diédral
d'ordre $2n$. Les images $R_i$ des sommets vérifient $R_i = R_{a_0(i)}$ et
forment un ensemble $J = \tilde J/\{1, a_0\}$ de $n$ points sur le cercle réel
$\Sigma_{\mathbb{R}} = \mathbb{P}(V)$, régulièrement répartis\footnote{Le
paragraphe est écrit en oblique, entremêlé d'une note marginale dont seuls
quelques mots se lisent ; on n'en retient que les formules.}.

\pagerange{46}{48}
\subsection*{Les pôles et la décomposition isotrope (pages 46 à 48)}

Les points de $\Sigma$ fixés par $D^{+}$ sont au nombre de deux, un dans
chacun des hémisphères que découpe $\Sigma_{\mathbb{R}}$. Ces hémisphères sont
indexés par $\underline{\omega}$ — l'hémisphère $\Sigma_\omega$ est celui dont
le bord, avec l'orientation de bord, a l'orientation $\omega$ — et les pôles
sont notés
\[
(15)\qquad P_\omega = \Sigma^{D^{+}} \cap \Sigma_\omega \qquad (\omega \in
\underline{\omega}).
\]
La page justifie (14), $\pi_0(\Sigma \smallsetminus \Sigma_{\mathbb{R}})
\simeq \underline{\omega}$, par ce qu'on lit « couvr. de Stokes » :
l'orientation induite sur le bord par l'orientation canonique de
$\Sigma$\footnote{Lecture incertaine du mot abrégé.}.
En algèbre linéaire, ce sont les droites propres de $D^{+}$ dans
$V_{\mathbb{C}}$. La forme quadratique s'étend en une forme bilinéaire complexe
non dégénérée, et $V_{\mathbb{C}}$ se décompose en ses deux droites isotropes,
\[
(16)\qquad V_{\mathbb{C}} = \bigoplus_{\omega \in \underline{\omega}}
V_{\mathbb{C}}^{\omega},
\]
qui sont aussi les droites propres du groupe $O^{+}(V)$ des rotations, et que
les éléments de $O^{-}(V)$ échangent. Une orientation $\omega$ définit le quart
de tour $u_\omega \in O^{+}(V)$, $u_\omega^2 = -1$, donc une structure complexe
$V_{(\omega)}$ sur $V$ ($i \cdot x = u_\omega x$) ; l'application
$\mathbb{C}$-linéaire $p_\omega : V_{\mathbb{C}} \to V_{(\omega)}$ qui prolonge
l'identité donne un isomorphisme $p = (p_\omega) : V_{\mathbb{C}} \simeq
\prod_\omega V_{(\omega)}$, et
\[
(19)\qquad V_{\mathbb{C}}^{\omega} = \ker p_{\omega'} = \{\, z \in
V_{\mathbb{C}} \mid u_\omega z = i z \,\} = \{\, x - i\,u_\omega x \mid x \in V
\,\}.
\]
On note $\varphi_\omega : V \xrightarrow{\sim} V_{\mathbb{C}}^{\omega}$,
$x \mapsto x - i\,u_\omega x$, l'isomorphisme $\mathbb{R}$-linéaire induit, et
encore $\varphi_\omega$ son prolongement $\mathbb{C}$-linéaire à
$V_{\mathbb{C}}$ ; on a $\operatorname{Re} \varphi_\omega(x) = x$ pour $x \in
V$, et c'est ce qui sert aux pages 35 à 39\footnote{Vérification : $u_\omega(x
- i u_\omega x) = u_\omega x + i x = i(x - i u_\omega x)$, et
$p_{\omega'}(x - i u_\omega x) = x - u_{\omega'} u_\omega x = 0$ puisque
$u_{\omega'} = u_\omega^{-1}$.}.

\pagerange{48}{51}
\subsection*{Le fibré $\mathcal{O}_\Sigma(1)$ et son torseur (pages 48 à 51)}

Soit $P\Sigma$ le $\mathbb{C}^{*}$-torseur associé à $\mathcal{O}_\Sigma(1)$.
En un point $x$ de $\Sigma$, correspondant à une droite $F$, les isomorphismes
canoniques (18)--(20) de la page — fibre de $\mathcal{O}(1)$ égale au quotient
par $F$, $F \otimes \mathcal{O}_F(1) \simeq \det V \simeq
\mathbb{R}(\underline{\omega})$, $\mathcal{O}_F(1) \simeq \check
F(\underline{\omega}) \simeq F^{\perp}$ — identifient $P\Sigma$ au plan privé de
l'origine,
\[
(21)\quad P\Sigma \simeq V^{*}, \qquad \text{au-dessus de} \qquad
(22)\quad \sigma : \Sigma \to \Sigma, \quad F \mapsto F^{\perp},
\]
avec, sur $V^{*}$, l'action opposée de $\mathbb{C}^{*}$\footnote{La page
réutilise les numéros (18), (19) et (21) ; on garde ses numéros. Elle écrit
« une droite $F$ de $V$ » sans préciser qu'il s'agit d'une droite complexe ;
le détail des torsions par $\underline{\omega}$ est en partie illisible, en
particulier dans la note marginale de la page 50.}. L'involution $\sigma$
fixe les deux pôles (les droites isotropes sont leurs propres orthogonales) et
fait sur le cercle réel un demi-tour (deux droites réelles orthogonales sont
diamétralement opposées sur $\mathbb{P}(V)$) : c'est le demi-tour autour de
l'axe des pôles, la « symétrie par rapport à l'axe polaire » de la page. Avec la
conjugaison complexe $\tau$, réflexion par le plan de l'équateur, on a
\[
\tau\sigma = \underline{a}, \qquad \text{l'antipodie de } \Sigma,
\]
puisque le demi-tour autour d'un axe composé avec la réflexion par le plan
orthogonal est la symétrie centrale.

On s'intéresse à
\[
(26)\qquad \Sigma^{*} = \Sigma \smallsetminus \{ R_i \mid i \in J \}, \qquad
(27)\qquad P\Sigma^{*} = P\Sigma|_{\Sigma^{*}} \simeq V^{**} = V^{*}
\smallsetminus \bigcup_{i \in J} \Delta_i,
\]
où les $\Delta_i$ sont les $n$ droites complexes des diagonales du
polygone. Pour $n = 3$, $\Sigma^{*}$ est la sphère privée de trois points,
c'est-à-dire $M_{0,4}$ (page 15), et $V^{**}$ le plan privé de trois droites.

\pagerange{52}{54}
\subsection*{Le programme, et les symétries (pages 52 à 54)}

Le but est de décrire des groupoïdes équivalents aux groupoïdes fondamentaux
\[
(28)\qquad \Pi_1 P\Sigma^{*} \simeq \Pi_1 V^{**}, \qquad \Pi_1 \Sigma^{*},
\]
et le morphisme (29) $\Pi_1 V^{**} \to \Pi_1 \Sigma^{*}$, par générateurs et
relations, avec des points de base stables par $D$ et par $\tau$. Sur la
situation agit le groupe $D \times \mathbb{C}^{*} \times \{1, \tau\}$, non
fidèlement : $(a_0, -1)$ agit trivialement. Son sous-groupe fini $D \times
\mu_m \times \{1, \tau\}$ est d'ordre $2m \cdot m \cdot 2 = 4m^2 = 16 n^2$, et
le quotient qui agit fidèlement, $(D \times \mu_m)/\{\pm1\} \times \{1,
\tau\}$, d'ordre $8n^2$ ($144$ et $72$ pour $n = 3$).

Sur $V_{\mathbb{C}} = \bigoplus_\omega V_{\mathbb{C}}^{\omega}$, le groupe
$\mu_m^{\underline{\omega}}$ agit coordonnée par coordonnée, et une rotation
$\zeta$ (comptée dans le sens $\omega$) agit par $\zeta$ sur
$V_{\mathbb{C}}^{\omega}$ et par $\zeta^{-1}$ sur $V_{\mathbb{C}}^{\omega'}$.
L'homomorphisme
\[
(34)\qquad D^{+} \times \mu_m \longrightarrow \mu_m^{\underline{\omega}},
\qquad (\zeta, \lambda) \longmapsto (\zeta\lambda,\ \zeta^{-1}\lambda),
\]
entre groupes d'ordre $m^2$, a pour noyau $\{(1,1), (a_0, -1)\}$, et pour
image, d'indice $2$, les couples tels que
\[
(35)\qquad (\lambda_\omega \lambda_{\omega'})^{n} = 1 :
\]
le produit des deux coordonnées est $\lambda^2$, un carré de $\mu_{2n}$, et
inversement tout couple dont le produit est un carré $\lambda^2$ s'écrit ainsi
avec $\zeta = \lambda_\omega/\lambda$. Les automorphismes de $V_{\mathbb{C}}$
qui stabilisent $V^{**}$, c'est-à-dire qui induisent sur $\Sigma$ un
automorphisme stabilisant $J$, proviennent de $D \times \mathbb{C}^{*}$ — le
stabilisateur de $n \geq 3$ points régulièrement répartis sur un cercle est le
groupe diédral d'ordre $2n$ — et ceux qui fixent chaque pôle, de $D^{+} \times
\mathbb{C}^{*}$, dont l'image dans $\mathbb{C}^{*\,\underline{\omega}}$ est
définie par $(\lambda_\omega \lambda_{\omega'}^{-1})^{n} = 1$ (37), condition
qui équivaut à (35) sur $\mu_m^{\underline{\omega}}$. La rédaction s'arrête
sur le choix de points de base stables par $D \times \{1, \tau\}$, et
« éventuellement » par le groupe d'ordre $72$ quand $n = 3$.

\pagerange{19}{32}
\section*{Générateurs et relations pour $n = 3$ (pages 19 à 32)}

Les pages 19 à 32 sont les calculs, pour l'hexagone ($n = 3$), que la
rédaction précédente encadre. Ils sont fragmentaires ; on en donne la
structure, et l'on dit ce qui s'y vérifie.

\emph{Les vecteurs.} Dans le plan $V = V(J)$ des vecteurs de somme nulle de
$\mathbb{R}^J$, $J = \{i, j, k\}$, la page 21 pose $\widetilde{R}_i = (e_j -
e_k) \wedge \omega_{jk}$, où $\omega_{jk}$ est l'orientation qui envoie $k$ sur
$j$ (la torsion par $\wedge\,\omega$ rend le vecteur indépendant de l'ordre
écrit), et $\widetilde{Q}^\omega_i = \widetilde{R}_{\omega^2 i} -
\widetilde{R}_{\omega i} = (e_j + e_k - 2e_i) \wedge \omega$ : les six vecteurs
$\pm\widetilde{R}_i$ sont les sommets d'un hexagone, les six
$\pm\widetilde{Q}_i$ les directions des milieux des côtés, et $i_\omega
\widetilde{Q}^\omega_i = \widetilde{R}_i$, $i_\omega \widetilde{R}_i =
-\widetilde{Q}^\omega_i$, où $i_\omega = u_\omega$ est le quart de tour. Le
dessin de la page 19 — un cylindre dont le cercle du haut porte
$V^{\omega}_{\mathbb{C}}$, celui du bas $V^{\omega'}_{\mathbb{C}}$, et
l'équateur $V$ — place les points de base : des points $\widetilde{P}^{\omega}_a$
au-dessus des pôles, $\widetilde{Q}_a$ et $\widetilde{R}_r$ sur l'équateur,
indexés par des ensembles $A$, $S$ et $R \simeq A \times \underline{\omega}
\simeq S \times \underline{\omega}$, sur lesquels agissent deux involutions
$\tau_0, \tau_1$ de produit $\omega = \tau_1\tau_0$\footnote{La page 19 est un
dessin d'une page entière, avec une légende : $r = (s,a)$, $r_1 = \tau_1 r$,
$r_2 = \omega^{-1}(r_1) = \tau_0\tau_1 r$, $r' = \tau_0 r$, $r'' = \tau_1\tau_0
r = \omega(r)$ ; la page 26 note $\sigma_0, \sigma_1$ ce que la page 19 note
$\tau_0, \tau_1$. L'ensemble à trois éléments $\{i,j,k\}$ est lu $J$ avec doute.
Les pages 23 (figures) et 21 reprennent un dessin d'avant ces pages, qu'on
n'a pas.}.

\emph{Les générateurs (page 26).} Des flèches « obliques »
\[
\tilde a_r : \widetilde{R}_r \to \widetilde{R}_{r'}, \qquad \tilde b_r :
\widetilde{Q}_a \to \widetilde{R}_r, \qquad \tilde c_r :
\widetilde{P}^{\omega}_a \to \widetilde{Q}_a,
\]
et des flèches « fibres » : les lacets $\lambda$ en chaque point de base (le
tour complet de la fibre $\mathbb{C}^{*}$), les demi-tours $\mu :
\widetilde{Q}_a \to \widetilde{Q}_a^{\vee}$, $\widetilde{R}_r \to
\widetilde{R}_r^{\vee}$ vers le point opposé $x^{\vee} = -x$, et les pas
$\tilde d$ entre points de base consécutifs au-dessus d'un même pôle.

\emph{Les relations.} Une relation « oblique » (1), $\tilde a_r \tilde b_r
\tilde c_r \tilde d_r = \tilde b_{r_1} \tilde c_{r_2}$, qui exprime qu'un
triangle de la figure se remplit ; des relations « fibres » : le produit de $m$
pas $\tilde d$ consécutifs est le lacet $\lambda$ (5), deux demi-tours
successifs font un tour, $\mu_{x^{\vee}}\mu_x = \lambda_x$ (6), (7) ; et des
relations « de transport » (8)--(10), qui disent que les flèches obliques
transportent le lacet de la fibre au lacet de la fibre\footnote{Les numéros des
relations fibres sont surchargés : (2) corrigé en (5), (5) en (6), (6) en (7).}.
Les relations fibres et de transport sont justes : le lacet de la fibre est
l'image du générateur de $\pi_1(\mathbb{C}^{*})$ par l'application orbite,
naturelle en le point, et il est donc central ; le demi-tour est le chemin
$t \mapsto e^{i\pi t} x$.

\emph{L'élimination (pages 28 à 32).} Posant $\widetilde{A}_r = \tilde
b_r^{-1} \tilde a_r^{-1} \tilde b_{r_1}$ et $\widetilde{B}_r = \tilde
a_{r''}^{-1} \tilde b_{r'} \tilde b_r^{-1}$, la page exprime, à l'aide de (1),
les demi-tours comme des produits
\[
\mu_{\widetilde{Q}_a} = \widetilde{A}_{\omega^{m} r''} \cdots
\widetilde{A}_{r''}, \qquad \mu_{\widetilde{R}_r} = \widetilde{B}_{\omega^{m} r}
\cdots \widetilde{B}_r,
\]
puis prend ces formules, et (5)--(7), pour \emph{définitions} des $\mu$ et des
$\lambda$. Il reste
\[
(1)\quad \tilde a_r \tilde b_r \tilde c_r \tilde d_r = \tilde b_{r_1} \tilde
c_{r_2}, \qquad (2')\quad \widetilde{A}_{\omega^{m} r''} \cdots
\widetilde{A}_{r''} = \widetilde{A}_{\omega^{m} r_1} \cdots \widetilde{A}_{r_1},
\]
et la page note que la relation (2''), $\tilde c^\omega_a(\lambda) = \tilde
c^\omega_a(\lambda')$ entre les lacets au-dessus des deux pôles, est, en
présence de (1), conséquence de (2') portée au carré. Le calcul de la page 32
sur deux paires de flèches $\mu, \mu' : Q \to Q^{\vee}$ et $\check\mu,
\check\mu'$ en retour — de $\check\mu\mu = \check\mu'\mu'$ on tire
$\check\mu'^{-1}\check\mu = \mu'\mu^{-1}$ — est juste. Les lignes qui
diraient si le système obtenu est une présentation complète de $\Pi_1
V^{**}$ sont illisibles (« ces deux relations \ldots{} équivalentes et $\pm$
triviales »), et le dossier n'énonce pas de présentation finale ; nous ne
l'affirmons pas\footnote{L'exposant des produits, $m$ ou $n$, change d'une
ligne à l'autre sur la page ; on écrit $m$, le nombre de pas d'un tour.
Les indices $0$ et $a$ des $\widetilde{Q}$ sont mêlés ; la marge de la page
28 parle de « 2 groupes de relations » une fois les $\mu$ éliminés.}. Pour
$n = 3$, le groupe à présenter est connu : $\pi_1(V^{**}) \simeq P_3 \simeq
F_2 \times \mathbb{Z}$ (page 17 ci-dessus).

\pagerange{35}{39}
\section*{Lemmes d'homotopies linéaires (pages 35 à 39)}

Ces lemmes servent à comparer, dans $\Pi_1 V^{**}$, un chemin et son image
par $\varphi_\omega$ (pages 46 à 48) : ils fournissent des homotopies
explicites qui évitent les droites ôtées. Pour $x, y \in V$, on pose
\[
\ell_{y,x}(t) = \tfrac{x+y}{2} + e(t)\,\tfrac{x-y}{2}, \quad e(t) = e^{i\pi
t}, \qquad L_{y,x}(t) = (1-t)x + t y,
\]
un demi-cercle complexe de $x$ à $y$ au-dessus du segment $[x,y]$, et le
segment lui-même. Dans ce qui suit, $\Delta$ est une droite réelle de $V$ et
$\Delta_{\mathbb{C}}$ son complexifié. Pour les droites ôtées de $V^{*}$,
$\Delta$ passe par l'origine.

\textbf{Lemme 2} (« trivial »). Soit $\ell$ un chemin dans $V$ ne rencontrant
pas $\Delta$, et $g(s,t) = (1-s)\varphi_\omega\ell(t) + s\,\ell(t)$. Alors $g$
évite $\Delta_{\mathbb{C}}$ : $\operatorname{Re} g(s,t) = \ell(t)$, puisque
$\operatorname{Re}\varphi_\omega(x) = x$, et $\operatorname{Re}
\Delta_{\mathbb{C}} = \Delta$. \emph{Corollaire} : si $\ell : x \to y$ est un
chemin de $V$ évitant les $\Delta_i$, le carré formé de $\varphi_\omega(\ell)$,
$\ell$ et des segments verticaux $L_{x, \varphi_\omega(x)}$, $L_{y,
\varphi_\omega(y)}$ commute dans $\Pi_1 V^{**}$.

\textbf{Lemme 3} (« $\pm$ trivial »). Si $\Delta \cap \{x, y\} = \emptyset$,
le demi-cercle $\ell_{y,x}$ évite $\Delta_{\mathbb{C}}$. En effet, $\ell_{y,x}(t)$
a pour partie imaginaire $\sin(\pi t)\,\frac{x-y}{2}$ et pour partie réelle un
point du segment $[x,y]$. S'il était dans $\Delta_{\mathbb{C}}$ pour $0 < t <
1$, la direction $x - y$ serait celle de $\Delta$, la droite $(xy)$ serait
parallèle à $\Delta$ et la rencontrerait, donc lui serait égale, et $x, y \in
\Delta$ ; pour $t = 0, 1$, le point est $x$ ou $y$\footnote{La page écrit
$\sin t$ et $\frac{y-x}{2}$ ; avec $e(t) = e^{i\pi t}$ de la page 35, c'est
$\sin(\pi t)\frac{x-y}{2}$. L'argument de la page, fragmentaire, est celui-ci.}.

\textbf{Lemme 4.} Soit $\ell_x(t) = e(t)\,x$, le demi-tour de $x$ à $-x$, et
$h(s,t) = e(t)\bigl((1-s)\varphi_\omega(x) + s x\bigr)$. Si $\Delta$ passe par
l'origine et $x \notin \Delta$, $h$ évite $\Delta_{\mathbb{C}}$ : comme
$\Delta_{\mathbb{C}}$ est stable par les homothéties complexes, $h(s,t) \in
\Delta_{\mathbb{C}}$ forcerait $(1-s)\varphi_\omega(x) + s x \in
\Delta_{\mathbb{C}}$, dont la partie réelle est $x$. \emph{Corollaire} : dans
$\Pi_1 V^{**}$, pour $x \in V^{**} \cap V$,

\begin{tikzcd}[column sep=huge, row sep=large]
\varphi_\omega(x) \arrow[r, "\ell_{\varphi_\omega(x)}"] \arrow[d, "L_{x,\varphi_\omega(x)}"'] & -\varphi_\omega(x) \arrow[d, "L_{-x,-\varphi_\omega(x)}"] \\
x \arrow[r, "\ell_x"] & -x
\end{tikzcd}

commute\footnote{La page énonce le lemme pour « toute droite $\Delta$ de $V$
telle que $x \notin \Delta$ » ; il faut que $\Delta$ passe par l'origine,
ce qui est le cas des droites ôtées. Elle étiquette la flèche verticale de
droite $L_{-x, \varphi_\omega(x)}$ ; c'est le segment de $-\varphi_\omega(x)$ à
$-x$, $h(s, 1) = -\bigl((1-s)\varphi_\omega(x) + sx\bigr)$.}.

\textbf{Lemme 1.} La page considère $f_{x,y}(s,t) = (1-s)\varphi_\omega
\ell_{y,x}(t) + s\,\ell_{y,x}(t)$ et affirme que, sous une hypothèse
d'orientation reliant le segment orienté $\overrightarrow{[x,y]}$, la droite
$\Delta$ et $\omega$, $f_{x,y}$ évite $\Delta_{\mathbb{C}}$ ; d'où, avec le
lemme 3, le carré commutatif
\[
\varphi_\omega(\ell_{y,x}) = \ell_{\varphi_\omega(y), \varphi_\omega(x)}
\quad\text{(linéarité de } \varphi_\omega\text{)}, \qquad
\ell_{y,x} \circ L_{x,\varphi_\omega(x)} = L_{y,\varphi_\omega(y)} \circ
\ell_{\varphi_\omega(y), \varphi_\omega(x)} .
\]
L'hypothèse est formulée et raturée plusieurs fois, et ne se lit pas. Elle est
indispensable : pour $x = (1,1)$, $y = (1,-1)$ dans une base orthonormée
directe pour $\omega$, et $\Delta$ l'axe des abscisses, on a $\ell_{y,x}(\frac12)
= (1, i)$, qui est dans $V_{\mathbb{C}}^{\omega'} = \ker\varphi_\omega$, donc
$f_{x,y}(0, \frac12) = 0 \in \Delta_{\mathbb{C}}$. Le lemme vaut donc pour une
seule des deux orientations du demi-cercle par rapport à $\omega$, et nous ne
reconstituons pas l'énoncé exact\footnote{Le contre-exemple est de cette
édition. Les premières lignes du lemme (« Soit $\overrightarrow{[x,y]}$ un
segment orienté \ldots{} tel que l'orientation de $\Delta$ induite \ldots{}
aille de $x$ vers $y$ ») et une note marginale oblique sont en partie
illisibles ; les flèches verticales du carré ont une tête à chaque bout sur la
page.}.

\pagerange{55}{77}
\section*{Feuillets de travail (pages 55 à 77)}

Ces pages ne forment pas une suite ; on les prend dans l'ordre, en n'énonçant
que ce qui se lit et se vérifie.

\pagerange{55}{65}
\subsection*{Cinq points sur la sphère et l'icosaèdre (pages 55 à 61, 63 et 65)}

Sous le titre « ${}_0\mathcal{T}_{0,5}$ et $S_0\mathcal{T}_{0,5}$ » (page 55),
les pages 57 à 61 recensent des configurations de points marqués et de
cercles sur la sphère — « $4$ sommets, deux faces », « $1$ cercle marqué, $6$
sommets, $4$ faces », des types I à IV — et comptent, pour chacune, les
« possibilités pour $K$ » : $60 = 12 \cdot 5$ (« pentagone $+$ sommet », «
pentagone $+$ arête »), $60 = 10 \times 3 \times 2$, $12$, $20$, $30$. Elles
nomment ces données par l'icosaèdre : un sommet et une arête incidente (type
$D_{01}$), un sommet et une face incidente ($D_{02}$), avec des flèches
d'oubli $D_{01} \to D_0$, $D_{01} \to D_1$, $D_{02} \to D_0$, $D_{02} \to D_2$,
$D_1 \to D_2$ — des drapeaux partiels de l'icosaèdre\footnote{« $12.5 = 60$ »
sous II se lit aussi « $66$ » ; le calcul donne $60$. Les croquis du bas de la
page 59 sont largement raturés.}.

La page 61 et le tableau de la page 65 donnent les stabilisateurs, dans le
groupe des rotations de l'icosaèdre ($\mathfrak{A}_5$, d'ordre $60$) puis dans
le groupe complet ($\mathfrak{A}_5 \times \mathbb{Z}/2$, « non orienté »), et
les nombres d'orbites correspondants :
\[
\begin{array}{l|cc|cc}
 & \text{stab.\ orienté} & \text{nombre} & \text{stab.\ non orienté} & \text{nombre}\\ \hline
\text{sommets} & \mathbb{Z}/5,\ \text{paires : } D_5 & 12,\ 6 & D_5 \times \mathbb{Z}/2 & 6\\
\text{arêtes} & \mathbb{Z}/2,\ \text{paires : } D_2 & 30,\ 15 & D_2 \times \mathbb{Z}/2 & 15\\
\text{faces} & \mathbb{Z}/3,\ \text{paires : } \mathfrak{S}_3 & 20,\ 10 & \mathfrak{S}_3 \times \mathbb{Z}/2 & 10
\end{array}
\]
Les nombres se vérifient : $60/5 = 12$, $60/10 = 6$, $60/2 = 30$, $60/4 = 15$,
$60/3 = 20$, $60/6 = 10$, et dans le groupe d'ordre $120$ les stabilisateurs
des paires antipodales sont deux fois plus gros\footnote{Le tableau de la page
65 a quatre colonnes de groupes et quatre colonnes de nombres, « NB des cas »
et « Nb de feuillets » ; on l'a recomposé selon ce qu'il compte. Une entrée de
la deuxième ligne pourrait se lire $\mathbb{Z}/4$. La feuille porte au verso un
imprimé daté du 31 mai 1979.}. La page 63 dessine l'icosaèdre, un patron
déplié aux arêtes marquées $a, \ldots, e$, et un « pentagone orienté ».
Le lien avec $M_{0,5}$ — sur lequel le groupe $\mathfrak{S}_5$ agit, et dont
$\mathfrak{A}_5$ est le sous-groupe d'indice $2$ — n'est pas écrit.

\pagerange{62}{62}
\subsection*{Points base à l'infini (page 62)}

« Avant de me lancer vers une exploitation des principes dégagés, il me faut
encore développer la formulation des \emph{points base à l'infini}. » Le point
de départ est la dimension $1$, un retour en dimension quelconque, pour un
diviseur à croisements normaux, étant annoncé. Soit $X$ le spectre d'un anneau
de valuation discrète hensélien, $x$ son point fermé, $k = k(x)$, $\mathcal{U}
= X \smallsetminus \{x\}$ son point générique, $T_x =
\mathbb{V}(\mathfrak{m}_x/\mathfrak{m}_x^2)$ l'espace tangent et $T_x^{*} = T_x
\smallsetminus \{0\}$. Il y a un morphisme canonique de groupoïdes
fondamentaux
\[
(55)\qquad \widehat{\Pi}_1(T_x^{*}) \longrightarrow \widehat{\Pi}_1(\mathcal{U}),
\]
et la page s'arrête sur « Alors on sait qu'on a ». Ce qu'on sait aujourd'hui :
en caractéristique résiduelle nulle, les deux membres sont des extensions de
$\mathrm{Gal}(\bar k/k)$ par $\widehat{\mathbb{Z}}(1)$, et le morphisme
canonique est une équivalence ; en général, il en est une sur les quotients
modérés. Un point de
$T_x^{*}$ — un vecteur tangent non nul — sert ainsi de point base pour le
point générique : c'est le \emph{point base tangentiel} que Deligne
formalisera en 1989\footnote{L'énoncé sur (55) et le nom sont de cette
édition ; la page ne dit pas comment (55) est défini et passe par le complété
$\widehat{\mathcal{U}}$, l'isomorphisme $\widehat{\Pi}_1(\widehat{\mathcal{U}})
\to \widehat{\Pi}_1(\mathcal{U})$ étant lu avec doute. Les numéros (53) à (55)
prolongent une numérotation qui n'est pas dans le dossier.}. C'est le
fondement, sur un corps quelconque, des points de base $\varepsilon$ du
dossier.

\pagerange{64}{66}
\subsection*{Le monogone et le bigône (pages 64 et 66)}

La page 64 compose $\varphi(z) = \frac{z-1}{z+1}$ ($1 \mapsto 0$, $-1 \mapsto
\infty$, $0 \mapsto -1$, $\infty \mapsto 1$), le carré, et $\psi(w) =
\frac{w}{w-1}$ ($0, 1, \infty \mapsto 0, \infty, 1$) :
\[
(26)\qquad g_1(z) = \psi\bigl(\varphi(z)^2\bigr) = -\frac14 \frac{(z-1)^2}{z},
\qquad
(27)\qquad g_2(z) = g_1(z^2) = -\frac14\Bigl(z - \frac1z\Bigr)^2,
\]
« équation du monogone » et « équation du \emph{bigône}, de merveille ! ».
Les calculs sont justes, et ce sont des fonctions de Belyi. $g_1$, de degré
$2$, a pour points critiques $z = 1$ (valeur $0$) et $z = -1$ (valeur $1$), et
$g_1^{-1}(\infty) = \{0, \infty\}$ : l'image réciproque du segment $[0,1]$ est
une boucle issue de $z = 1$, passant par $z = -1$, qui sépare la sphère en
deux faces centrées en $0$ et $\infty$ — le monogone, carte à un sommet, une
arête, deux faces. $g_2$, de degré $4$, a deux sommets $\pm 1$, deux milieux
d'arêtes $\pm i$ et deux faces centrées en $0$ et $\infty$ : le bigône. La
page 66 vérifie la valeur $g_1(i) = -\frac{1}{4i}(i-1)^2 = \frac12$, ($(i-1)^2 =
-2i$), et tabule les images de $0, \pm 1, \infty$ par $X_2 \xrightarrow{z^2}
X_1 \to X_0$\footnote{La première ligne de la page 64 suppose une
normalisation antérieure (« les étaient $0, \infty$ au lieu de $0, 1$ ») et
renvoie à une formule (25) absente du dossier. La page 66, écrite dans les
deux sens, porte aussi un quotient griffonné dont le signe et le dénominateur
ne se vérifient pas tels qu'on les lit ; on ne le reproduit pas.}.

La moitié tête-bêche de la page 66 nomme les deux catégories : « Cartes,
$\mathcal{C}_2$-ens » et « Cartes orientées, $\mathcal{C}_2^0$-ens », avec
\[
\mathcal{C}_2 = \langle \sigma_0, \sigma_1, \sigma_\infty \mid \sigma_0^2 =
\sigma_1^2 = \sigma_\infty^2 = (\sigma_0\sigma_\infty)^2 = 1 \rangle, \qquad
\mathcal{C}_2^0 = \langle \rho_0, \rho_1, \rho_\infty \mid \rho_\infty\rho_1\rho_0
= \rho_1^2 = 1 \rangle .
\]
C'est le groupe cartographique des cartes, celui de la page 9 du dossier 141
(qui l'écrit avec les indices $0, 1, 2$) et de la page 22 du dossier 145, et
son sous-groupe orienté $\mathbb{Z} * \mathbb{Z}/2$ ; la relation
$\rho_\infty\rho_1\rho_0 = 1$ est dans l'ordre de la seconde rédaction du
dossier 145. Une carte orientée finie est un ensemble fini sur lequel agit
$\mathcal{C}_2^0$, et les cartes de $g_1$ et $g_2$ en sont les plus
petits exemples\footnote{Les présentations sont sur la page ; leur
identification au groupe cartographique est de cette édition.}.

\pagerange{67}{67}
\subsection*{L'espace $\widehat{M}_{0,5}$ (page 67)}

$M_{0,5} \subset \widehat{M}'_{0,5} \subset \widehat{M}_{0,5}$, où
$\widehat{M}'_{0,5}$ est l'ouvert des courbes stables de type $(0,5)$ qui
restent stables quand on oublie le cinquième point. On a
\[
\widehat{M}'_{0,5} \simeq \widehat{\mathcal{X}}^{*}_{0,4} =
\widehat{\mathcal{X}}_{0,4} \smallsetminus (\tilde\Delta_0 \cup \tilde\Delta_1
\cup \tilde\Delta_\infty \cup \tilde\Delta),
\]
où $\widehat{\mathcal{X}}_{0,4} \to \widehat{M}_{0,4}$ est la courbe
universelle, éclatée de $\mathbb{P}^1 \times \mathbb{P}^1$ en les trois points
$(0,0), (1,1), (\infty,\infty)$ de la diagonale, et les $\tilde\Delta$ les
transformées strictes des quatre sections. C'est juste : $\overline{M}_{0,5}$
est la courbe universelle sur $\overline{M}_{0,4}$, et l'on retire les
diviseurs où le cinquième point rencontre l'un des quatre autres. La page se
demande ce qui manque — les courbes qui ne restent pas stables — et en dessine
les configurations.

\pagerange{70}{73}
\subsection*{Les fibrés $L_i$ et le groupe de Picard équivariant (pages 70 à
73)}

Sur $M_{0,I}$, les fibrés en droites $L_i$ (attachés aux points marqués) et
$L = \bigotimes_i L_i$ vérifient $L_i \simeq L_K \otimes p^{*}_{K
\smallsetminus\{i\}, K} L_{K \smallsetminus \{i\}}$. Un tableau cherche, pour
$n = \operatorname{card} I = 4, 5, 6$, le plus petit exposant qui trivialise
$L$, resp.\ $L_i$, de façon compatible à $\mathfrak{S}_I$ : $3$ et $6$ pour
$n = 4$, des valeurs suivies de points d'interrogation pour $n = 5, 6$, en
regard de $(n-1)(n-2) = 6, 12, 20$\footnote{Les deux dernières lignes du
tableau, « exposant pour trivialiser $H^1(M_{0,I}, \mathfrak{S}_I;
\mathbb{G}_m)$ » (biffée) et « dividende (?) », sont vides. La page 72 est
écrite sous un sujet d'examen de juin 1978 sur les classes de conjugaison de
$\mathfrak{A}_5$ et l'icosaèdre.}. La page 72 en donne la source : des
sections $\xi_A$ indexées par les parties $A$ à trois éléments, et leur
produit, section de $L^{\otimes\nu}$, $\nu = (n-1)(n-2)/2$, d'où « $L^{\otimes
6} \simeq 1$ » pour $n = 4$ et $n = 5$. Ces affirmations ne sont pas
démontrées sur les pages et ne sont pas vérifiées ici.

Le cadre est la cohomologie équivariante. Comme $M_{0,K}$ est le complémentaire
d'un arrangement d'hyperplans dans un espace affine, $\mathrm{Pic}(M_{0,K}) = 0$
(« $H^1(M_{0,5}; \mathbb{G}_m) \simeq 0$ »), et la suite spectrale de
$\mathfrak{S}_K$ donne $H^1(M_{0,K}, \mathfrak{S}_K; \mathbb{G}_m) \simeq
H^1(\mathfrak{S}_K, H^0(M_{0,K}, \mathbb{G}_m))$. De la suite exacte $0 \to
\mathbb{G}_m \to H^0(M_{0,K}, \mathbb{G}_m) \to U' \to 0$, où $U'$ est le
groupe des unités modulo les constantes (noté $(\mathbb{Z}^{\mathbb{P}_3(K)})'$),
on tire
\[
0 \to \underbrace{H^1(\mathfrak{S}_K, \mathbb{C}^{*})}_{\pm 1} \to
H^1(M_{0,K}, \mathfrak{S}_K; \mathbb{G}_m) \to H^1(\mathfrak{S}_K, U') \to
H^2(\mathfrak{S}_K, \mathbb{C}^{*}),
\]
exacte à gauche parce que $U'$ n'a pas d'invariant non nul sous
$\mathfrak{S}_K$\footnote{Cette justification est de cette édition :
$U' \otimes \mathbb{Q}$ est la représentation de $\mathfrak{S}_K$ indexée par
la partition $(n-2, 2)$, sans invariant. La page écrit la dernière flèche
verticalement, et indexe une fois $\mathfrak{S}_5$ pour $\mathfrak{S}_K$. Elle
porte aussi le plongement $M_{0,I} \hookrightarrow \prod_i \Sigma_{I
\smallsetminus\{i\}}$ par les applications d'oubli, juste, et $\Omega^1 \simeq
\mathcal{O}(-2)$ sur $\mathbb{P}^1$.}.

\pagerange{74}{77}
\subsection*{Groupes des polyèdres réguliers (pages 74, 76 et 77)}

La page 74 commence au milieu d'un énoncé, avec deux automorphismes encadrés
(41), (42) d'un groupe à générateurs $\rho_0, \rho_1, \rho_\infty$, $\tau_0,
\tau_\infty$, dont une image est laissée en blanc ; puis une liste de groupes
$\pi_{D_n}$, $\pi_{\mathfrak{S}_4^{+}}$, $\pi_{\mathfrak{S}_4}$,
$\pi_{\mathfrak{S}_5^{+}}$, et $\pi_X$ pour « $X$ une carte : on enlève les
sommets », avec $\pi_{\mathrm{icos}}$, $\pi_{\text{dodéc}}$,
$\pi_{\text{tétr}}$, $\pi_{\mathrm{cube}}$, $\pi_{\mathrm{oct}}$. Ce sont
vraisemblablement les groupes fondamentaux de la sphère privée des sommets
d'un polyèdre régulier, sur lesquels agissent les groupes de rotations
$D_n$, $\mathfrak{A}_4 = \mathfrak{S}_4^{+}$, $\mathfrak{S}_4$, $\mathfrak{A}_5 =
\mathfrak{S}_5^{+}$\footnote{Lecture de cette édition. Les groupes orbifolds
correspondants sont les groupes de triangles $(2,2,n)$, $(2,3,3)$, $(2,3,4)$,
$(2,3,5)$.}. La page 76 recense le dodécaèdre : $120$ repères (drapeaux),
$12$ faces, $30$ arêtes, $20$ sommets, et $2 \times (6 + 15 + 10) = 62 = 12 + 30
+ 20$, le nombre total de cellules\footnote{La page écrit « bidodécagones » ;
le nombre $12$ des faces est surchargé.}. La page 77, des croquis de cartes
et de rubans sans texte, est dessinée sur une circulaire du 12 février 1982.

\pagerange{78}{98}
\section*{« Groupoïdes » (pages 78 à 98)}

\pagerange{80}{80}
\subsection*{Présentation d'une extension de groupes (page 80)}

\textbf{Lemme 1.} Soit $1 \to L \to \tilde\pi \to \pi \to 1$ une suite exacte
de groupes, $(x_i)_{i \in I}$ des générateurs de $\pi$ de relations
définissantes $R_\alpha((x_i)) = 1$ ($\alpha \in A$), $(a_j)_{j \in J}$ des
générateurs de $L$ de relations définissantes $S_\beta((a_j)) = 1$ ($\beta \in
B$). On relève les $x_i$ en des $\tilde x_i \in \tilde\pi$, et l'on choisit des
mots $M_\alpha$, $M_{ij}$ en les $a_j$ tels que
\[
(1)\quad R_\alpha((\tilde x_i)) = M_\alpha((a_j)), \qquad (2)\quad \tilde
x_{i}\, a_{j}\, \tilde x_{i}^{-1} = M_{ij}((a_j)).
\]
Alors (1), (2) et (3) $S_\beta((a_j)) = 1$ sont des relations définissantes de
$\tilde\pi$ en les $\tilde x_i$, $a_j$. Si $L$ est central, (2) devient
$[\tilde x_i, a_j] = 1$ (2').

C'est le lemme classique de présentation d'une extension, et il est juste tel
quel, sans relation de conjugaison par les $\tilde x_i^{-1}$. Soit $\Pi$ le
groupe présenté et $N$ le sous-groupe engendré par les $a_j$. Par (3), $a_j
\mapsto a_j$ définit $L \to N$, et sa composée avec $\Pi \to \tilde\pi$ est
l'identité de $L$ : $N \simeq L$. Par (2), la conjugaison par $\tilde x_i$
envoie $N$ dans $N$, et elle correspond à la conjugaison par $\tilde x_i$ dans
$\tilde\pi$, qui est un automorphisme de $L$ ; elle est donc surjective, et
$N$ est distingué. Enfin $\Pi/N$ a pour présentation $\langle x_i \mid
R_\alpha \rangle$, c'est $\pi$, et $\Pi \to \tilde\pi$ est un
isomorphisme\footnote{La démonstration est de cette édition ; la page n'en a
pas. Les trois annotations marginales, « relèvements » des $R_\alpha$,
relations « d'action », relations de $L$, sont lues avec doute.}.

\pagerange{81}{86}
\subsection*{Présentation d'un groupoïde au-dessus d'un groupoïde (pages 81 à
86)}

La même chose pour un morphisme de groupoïdes $p : \widetilde{\mathcal{T}}
\to \mathcal{T}$ (la page 81, barrée, est un premier départ de la page 82).
On se donne des flèches $\tilde u_i$ de $\widetilde{\mathcal{T}}$ dont les
projections $u_i$ engendrent $\mathcal{T}$, avec des relations définissantes
de $\mathcal{T}$ écrites comme des lacets composables $u'_{i'_1} \cdots
u'_{i'_n} = 1$, où $u'_{(i,\varepsilon)} = u_i^{\varepsilon}$ (8). Les relevés
$\tilde u'$ ne sont pas composables, mais les extrémités consécutives sont
dans une même fibre ; on suppose donnée une famille de flèches fibres $a_j$
engendrant chaque groupoïde fibre $\widetilde{\mathcal{T}}_s$, et les fibres
connexes, de sorte que des mots $M_{\alpha,\nu}((a_j))$ relient ces extrémités
(11). On écrit alors :
\begin{itemize}
\item les relations « remontées » (12) : $M_{\alpha,1}\, \tilde u'_{i'_1}\,
M_{\alpha,2}\, \tilde u'_{i'_2} \cdots M_{\alpha,n}\, \tilde u'_{i'_n} = 1$ ;
\item des relations (13) $N'_\beta((a_j)) = 1$ qui définissent chaque fibre ;
\item pour chaque lacet fibre $a_{j_0}$ et chaque $\tilde u_i$ de même
origine, une relation de transport (14) $\tilde u_i(a_{j_0}) =
M''_{i,j_0}((a_j))$, le transporté étant un lacet de la fibre de
l'extrémité.
\end{itemize}
« Cela posé, je dis que (12) (13) (14) est un système de relations
définissant » $\widetilde{\mathcal{T}}$\footnote{La page écrit $\mathcal{T}$,
sans tilde ; c'est $\widetilde{\mathcal{T}}$ qui est présenté. Elle suppose
aussi que les lacets $a_j$ en chaque objet engendrent son groupe
d'automorphismes dans la fibre.}. La page s'arrête là, sans démonstration ;
c'est l'analogue en groupoïdes du lemme de la page 80, et c'est le cadre
exact de l'élimination des pages 26 à 32, où les $\lambda$, $\mu$, $\tilde d$
sont des flèches fibres et (8)--(10) des relations de transport.

\pagerange{88}{94}
\subsection*{$\pi$-fibrations de groupoïdes (pages 88 à 94)}

Sous le titre « Généralités sur les groupoïdes ». Soit $f : C' \to C$ un
morphisme de groupoïdes et $\pi$ un groupe. Une \emph{$\pi$-fibration} est un
$f$ surjectif sur les objets et plein, dont les fibres sont des groupoïdes
connexes munis d'un isomorphisme naturel de $\pi$ sur leurs groupes
d'automorphismes : on a donc un homomorphisme $\pi \to
\mathrm{End}(\mathrm{id}_{C'})$ — $\pi$ agit naturellement sur chaque objet
de $C'$ —, une suite exacte
\[
1 \to \pi \to \mathrm{Aut}_{C'}(x) \to \mathrm{Aut}_{C}(f(x)) \to 1,
\]
et $\pi_0(f)$ bijectif. $C'$ est une \emph{extension centrale} de $C$ par
$\pi$ (« l'extension est ici nécessairement centrale ») ; pour une extension
non centrale il faudrait remplacer $\pi$ par un système local de groupes sur
$C$, comme le dit la marge. En termes d'aujourd'hui, $C'$ est une gerbe liée
par $\pi$ au-dessus de $C$\footnote{La page écrit
$\mathrm{Aut}_{C'}(f(x))$ pour le dernier terme : c'est
$\mathrm{Aut}_C(f(x))$. La surjectivité à droite demande que $f$ soit plein,
ce que la page ajoute dans l'interligne (« et (sur les flèches) »). « Gerbe
liée » est de cette édition.}.

\emph{L'exemple (page 88).} Soit $U$ un groupe topologique connexe et
localement simplement connexe, $\pi = \pi_1(U, 1)$ (abélien), $P \to T$ un
$U$-torseur. La suite exacte d'homotopie $\pi \to \pi_1(P_i) \to \pi_1(T_i) \to
1$ fait de $\Pi_1(P) \to \Pi_1(T)$ une $\pi$-fibration quand les $\pi \to
\pi_1(P_i)$ sont injectifs, une $\pi/N$-fibration si leur noyau $N$ ne dépend
pas de la composante. Soit $\mu$ un groupe, $\mu \to U$ un homomorphisme,
$\varepsilon$ un ensemble à action de $\mu$ et $\varepsilon \to P$ une
application équivariante\footnote{Les pages 89 à 96 notent cet ensemble $E$,
d'une lettre appuyée qu'elles distinguent du $\mathcal{E}$ de l'extension ; on
écrit $\varepsilon$, comme aux pages 4 et 97.} ; soit $C'$ le groupoïde induit
par $\Pi_1(P)$ sur $\varepsilon$, et $C$ celui induit par $\Pi_1(T)$ sur
$\varepsilon_0 = \varepsilon/\mu$. Le carré
\[
C' \to \Pi_1(P), \qquad C \to \Pi_1(T)
\]
commute, ses flèches horizontales sont des équivalences quand $\varepsilon$
rencontre toutes les composantes, et $\mu$ agit sur $C'$ par transport de
structure. Soit $\mathcal{E}$ l'image inverse de $0 \to \pi \to \widetilde{U}
\to U \to 0$ par $\mu \to U$ (13). Si $\mu$ agit librement sur $\varepsilon$
et que $x, y$ sont dans une même fibre $C'_\xi$, on a un isomorphisme
canonique de $\pi$-torseurs
\[
(14)\qquad \mathrm{Isom}_{C'_\xi}(x, y) \simeq \mathcal{E}_{y - x},
\]
compatible aux compositions, c'est-à-dire $C'_\xi \simeq
\mathcal{E}(\varepsilon_\xi)$ (15) : c'est, en famille, la formule (9) des
pages 4 et 5.

\emph{Les $(\pi, \mu)$-fibrations (pages 92 à 94).} D'où la définition : une
$(\pi, \mu)$-fibration subordonnée à une extension centrale $1 \to \pi \to
\mathcal{E} \to \mu \to 1$ est une $\pi$-fibration $f : C' \to C$ munie d'une
action de $\mu$ sur $C'$ compatible aux opérations internes de $\pi$, telle que
$f$ soit $\mu$-équivariante pour l'action triviale sur $C$, que $\mu$ agisse
librement sur $\mathrm{Ob}\,C'$ avec $\mathrm{Ob}\,C'/\mu \simeq
\mathrm{Ob}\,C$, et qu'une condition sur les fibres, « à examiner », soit
remplie. Pour la formuler, la page analyse une fibre : un $\pi$-groupoïde
connexe $\Phi$, avec une action de $\mu$ dont l'ensemble des objets est un
$\mu$-torseur. Les $\pi$-torseurs $\mathrm{Hom}(x, y)$ sont identifiés entre
eux par l'action de $\mu$ quand la « différence » $u$ de $x$ et $y$ est fixée,
d'où des $H_u$, et la composition donne $H_u \wedge H_v \simeq H_{uv}$ avec
associativité et unité : une extension $\mathcal{E}_\Phi$ de $\mu$ par $\pi$.
Inversement $\mathcal{E}$ et un $\mu$-torseur $\varepsilon$ donnent
$\mathcal{E}(\varepsilon)$. D'où la correspondance
\[
(17)\qquad
\left\{\begin{array}{l}
\pi\text{-groupoïdes connexes } \Phi \text{ avec action de } \mu\\
\text{dont l'ensemble des objets est un } \mu\text{-torseur}
\end{array}\right\}
\longleftrightarrow
\{ \text{extensions de } \mu \text{ par } \pi \},
\]
que la page dit être une équivalence de catégories. Elle l'est si l'on munit
$\Phi$ d'un objet de base et qu'on demande aux morphismes de le
respecter ; sans cela, le membre de gauche a davantage de morphismes (le
choix de l'image d'un objet), et l'on n'a qu'une bijection entre classes
d'isomorphisme\footnote{Cette précision est de cette édition ; la page ne dit
pas quels sont les morphismes. Elle note qu'on n'a « pas utilisé la
commutativité de $\mu$ ».}. Une $\mathcal{E}(\pi, \mu)$-structure sur $C'$
au-dessus de $C$ est alors la donnée, pour chaque $\xi \in \mathrm{Ob}\,C$,
d'un isomorphisme $\mathcal{E}_{C'_\xi} \simeq \mathcal{E}$, avec une condition
de passage d'un $\xi$ à l'autre le long des flèches de $C$, « je présume »,
que le dossier ne formule pas.

\pagerange{95}{98}
\subsection*{Le quotient par $\mu$ (pages 95 à 98)}

Pour préciser cette condition, la page revient à l'exemple, avec $\mu \subset
U$ discret et distingué (18). $P_0 = P/\mu$ est un torseur sous $U_0 = U/\mu$,
le revêtement universel de $U_0$ est encore $\widetilde{U}$, et
\[
(20)\qquad \mathcal{E} \simeq \pi_1(U_0, 1),
\]
ce qui montre au passage que $\mathcal{E}$ est commutative, comme tout groupe
fondamental d'un groupe topologique. $\Pi_1(P_0)$ porte donc une
$\mathcal{E}$-structure (21), et les deux actions de $\mathcal{E}$ sur
$\mathrm{Hom}(x_0, y_0)$ se correspondent par $\ell\cdot\gamma_{x_0} =
{}^{\ell}\gamma_{y_0} \cdot \ell$. L'image inverse de $\Pi_1(P_0)$ par $P \to
P_0$ est $\Pi_1(P) \wedge^{\pi} \mathcal{E}$, et celle par $\varepsilon \to P
\to P_0$ est $C' \wedge^{\pi} \mathcal{E}$. Comme $\mu$ agit sur
$\mathcal{E}$ (trivialement ici sur $\pi$), il agit sur $C' \wedge^\pi
\mathcal{E}$, librement sur les objets, et le quotient est un
$\mathcal{E}$-groupoïde sur $\varepsilon_0$ — « plus de doute, c'est $C_0$ ! »,
le groupoïde induit par $\Pi_1(P_0)$ sur $\varepsilon_0$, ainsi construit
algébriquement à partir de $C'$ et de ses structures. La construction inverse
(« Inversement, partons d'un $\mathcal{E}$-groupoïde $C_0$ ») est commencée et
laissée.

La page 98, de notes, en donne les données : 1°) un $\mathcal{E}$-groupoïde
central $C_0$ sur $\varepsilon_0$ ; 2°) un revêtement principal $\varepsilon \to
\varepsilon_0$ de groupe $\mu$, donné par un système local sur $C_0$, c'est-à-dire
un homomorphisme $\pi_1(C_0, x_0) \to \mu$ qui prolonge $\mathcal{E} \to \mu$ ;
au niveau des groupes, $1 \to \mathcal{E} \to \pi_1(P_0) \to \pi_1(T) \to 1$ et
$1 \to \pi \to \ker\varphi \to \pi_1(T) \to 1$ pour $\varphi : \pi_1(P_0) \to
\mu$. Elle esquisse enfin la généralisation à un torseur sous un faisceau de
groupes topologiques $U_T$, de système local de groupes fondamentaux
$\Pi_T$, avec un sous-faisceau $\gamma_T$ et $0 \to \Pi_T \to \mathcal{E}_T \to
\gamma_T \to 1$\footnote{La page 97 écrit « sur le quotient $\varepsilon =
\varepsilon'/\mu$ », la lettre petite et ambiguë ; c'est $\varepsilon_0 =
\varepsilon/\mu$. Les marges des pages 89 et 95 sont en partie illisibles.}.

Ces pages sont l'outil du dossier sous sa forme générale : un groupoïde
fondamental induit sur un ensemble fini de points de base stable par un
groupe d'opérateurs $\mu$ — les $\varepsilon$ des pages 2 à 14 et 111 à 124,
avec $\mu = \{\pm1\}^{\ldots}$ et $U = \mathbb{U}^{\ldots}$ — se reconstruit
algébriquement à partir d'une extension et d'un torseur.

\pagerange{99}{124}
\section*{$S_0\mathcal{T}_G$ en général et les morphismes de généralisation
(pages 99 à 124)}

\pagerange{99}{100}
\subsection*{Graphes à sommets de genre $0$ et d'ordre au plus $4$ (pages 99 et
100)}

Pour un graphe MD dont tous les sommets sont de genre $0$ et d'ordre $3$ ou
$4$, la page 99 écrit $S_0\mathcal{T}_G$ comme le produit des
$S_0\mathcal{T}_{0, \hat{\vec A}_s}$, groupoïdes à opérateurs
$(\mathbb{Z}^{\hat{\vec A}_s}, \mu_2^{\hat{\vec A}_s})$, poussé le long de
$(\mathbb{Z}^{\hat{\vec A}}, \mu_2^{\hat{\vec A}}) \to (\mathbb{Z}^{\hat A},
\mu_2^{\hat A})$ : un groupoïde à opérateurs $(\mathbb{Z}^{\hat A},
\mu_2^{\hat A})$, sur les objets duquel $\mu_2^{\hat A}$ agit librement, et
dont le quotient par tous les opérateurs est
\[
R_0\mathcal{T}_G = \prod_{s \in S} \mathcal{T}_{0, \hat{\vec A}_s},
\]
où seuls comptent les sommets d'ordre $4$, ceux d'ordre $3$ donnant le
groupoïde ponctuel. Pour un graphe quelconque, de genre virtuel nul, la page
annonce $S_0\mathcal{T}_G$ comme limite inductive filtrante de tels groupoïdes,
indexée par un ensemble ordonné d'arbres finis sans sommet d'ordre $2$, et
s'arrête sur des points de suspension\footnote{Le premier membre est écrit sur
un symbole biffé, et l'exposant du produit contracté porte des signes en partie
biffés ; on ne restitue que la structure. La page 100 ne porte que le titre
« Hom de généralisation $SP_G \to SP_{G^\natural}$ ».}.

\pagerange{102}{103}
\subsection*{Réductions (pages 102 et 103)}

Le fibré $SP_G$ et tous ses quotients $P_{G,K}$ se déduisent d'un fibré plus
fin $SP_{G^b}$, de groupe $\mathbb{G}_m^{\vec A}$, par $\mathbb{G}_m^{\vec A}
\to \mathbb{G}_m^{A}$ ; donc $S\mathcal{T}_G$ se déduit de
$S\mathcal{T}_{G^b}$ par $\mathbb{Z}^{\vec A} \to \mathbb{Z}^A$. Comme
$\mathcal{T}_G$ est une extension de $R\mathcal{T}_G$ par $\mathbb{Z}^A$,
connue dès qu'on connaît les extensions par $\mathbb{Z}$ attachées à chaque
arc, tout se ramène à construire les $S\mathcal{T}_{g,I} \simeq
\mathcal{T}^{!}_{g,\nu}$, extensions de $\mathcal{T}_{g,I}$ par $\mathbb{Z}^I$
— ou les extensions par $\mathbb{Z}$ attachées à chaque $i \in I$\footnote{La
page qualifie $\mathcal{T}_G$ de « “cat.” de groupoïdes $R\mathcal{T}_G$ par
$\mathbb{Z}^A$ » ; « ext. » irait mieux avec la suite.}. De même, les
morphismes de généralisation $S\mathcal{T}_G \to S\mathcal{T}_{G'}$ se
ramènent au cas où $G' = G^{\natural}$ est la généralisation maximale de $G$
(toutes les arêtes liées contractées : la courbe lisse).

Pour une contraction $G \to G' = G/C$, $C \subset A$ (la section « Graphes
MD » ci-dessous), la page 103 dessine les tours de fibrés
\[
SP_G \xrightarrow{\ \mathbb{G}_m^{I}\ } P_G \xrightarrow{\ \mathbb{G}_m^{A'}\ }
P_{G,C} \xrightarrow{\ \mathbb{G}_m^{C}\ } RP_G = M_G,
\]
où $A' = A \smallsetminus C$ sont les arêtes de $G'$, et leurs images vers
celles de $G'$ et de $G^{\natural}$, les carrés étant cartésiens ; les flèches
horizontales ne sont définies que dans la catégorie homotopique, et commutent
aux groupes d'opérateurs. Par passage aux groupoïdes fondamentaux :
\[
S\mathcal{T}_G \xrightarrow{\ \mathbb{Z}^{I}\ } \mathcal{T}_G
\xrightarrow{\ \mathbb{Z}^{A'}\ } \mathcal{T}_{G,C} \xrightarrow{\
\mathbb{Z}^{C}\ } R\mathcal{T}_G, \qquad S\mathcal{T}_G \to S\mathcal{T}_{G'}
\to S\mathcal{T}_{G^\natural},
\]
avec $P_{G,K} = \prod_{a \in K} P_{G,a}$, $SP_G = P_{G,\hat A}$, $P_G =
P_{G,A}$, $SRP_G = P_{G,I}$, $RP_G = P_{G,\emptyset}$, et le carré cartésien
$SP_G \to SRP_G$, $P_G \to RP_G$ (groupes $\mathbb{G}_m^{A}$ horizontalement,
$\mathbb{G}_m^{I}$ verticalement)\footnote{Un premier état du grand diagramme,
en tête de la page 103, est barré en zigzag. La formule
$\mathcal{T}_{G,C} = \Pi_1 P_{G,\ldots}$ est coupée au bord de la feuille. La
page écrit ces fibrés avec $\mathbb{G}_m$ ; les pages 105 et suivantes, en
topologie, avec $\mathbb{U}$, ce qui ne change pas les groupoïdes
fondamentaux.}.

\pagerange{105}{111}
\subsection*{La chirurgie holomorphe (pages 105 à 111)}

Soit $G \to G' = G/C$ une contraction, et $X$ une courbe sur $\mathbb{C}$ de
type $G$ ; $C \subset A = \mathrm{Sing}(X)$, $B = A \smallsetminus C$. Soit
$\widetilde{X}$ la normalisée de $X$ en les points de $C$, et $\widetilde{C}$
l'image inverse de $C$. On choisit autour de chaque $x \in \widetilde{C}$ un
disque fermé $D_x$ à bord analytique réel (ou seulement $C^1$), ces disques
deux à deux disjoints et évitant les autres nœuds. L'espace de ces choix
devrait être contractible, c'est-à-dire que l'espace $M^C_G$ des courbes de
type $G$ munies de disques devrait se projeter sur $M_G$ par une équivalence
d'homotopie\footnote{« Il doit être vrai que » : la page ne le démontre pas,
et nous non plus ; c'est le seul point non démontré de la construction.}.
Alors $X_1 = \widetilde{X} \smallsetminus \bigcup D_x^{\circ}$ est une
surface holomorphe à bord, et chaque cercle $\partial D_x$ est un
$\mathbb{U}$-torseur, d'orientation opposée à celle de bord de $X_1$ ; il
s'identifie au torseur $T_x = (T_x\widetilde{X} \smallsetminus
0)/\mathbb{R}^{*}_{+}$ des directions tangentes, indépendamment du choix du
disque. On récupère $X$ en recollant à $X_1$ les cônes sur ces cercles.

Pour un nœud $x \in C$, de branches $x', x''$, l'espace des
anti-isomorphismes de torseurs
\[
\mathrm{Isom}^{-}_{\mathbb{U}}(T_{x'}, T_{x''}) \simeq
\mathrm{Isom}_{\mathbb{U}}(T_{x'}^{-1}, T_{x''}) \simeq T_{x'} \wedge T_{x''} =
T_x
\]
est un $\mathbb{U}$-torseur, et $P_C = \prod_{x \in C} T_x$ un
$\mathbb{U}^{C}$-torseur. Pour $t = (t_x) \in P_C$, on obtient une courbe
$X'_t$ de type $G'$ en recollant, pour chaque $x \in C$, les deux bords
$\partial D_{x'}$ et $\partial D_{x''}$ de $X_1$ selon l'anti-isomorphisme
$t_x$. C'est la construction qu'on appelle aujourd'hui l'\emph{ouverture des
nœuds} par recollement (« plumbing »), et $T_{x'} \wedge T_{x''}$ est le
torseur unitaire de $T_{x'} \otimes T_{x''}$, la droite où vit le paramètre de
lissage du nœud\footnote{Les noms et la dernière identification sont de cette
édition.}. Quand $(X, (D_x))$ varie, les $P_C$ forment un
$\mathbb{U}^C$-torseur $\widetilde{M}^C_G$ sur $M^C_G$, et l'on obtient un
morphisme $\widetilde{M}^C_G \to M_{G'}$, avec
\[
1 \to \mathbb{Z}^C \to \pi_1(\widetilde{M}^C_G) \to \pi_1(M^C_G) \simeq
\pi_1(M_G) \to 1, \qquad \pi_1(\widetilde{M}^C_G) \to \pi_1(M_{G'}).
\]

Pour relever aux fibrés tangents, soit $P_B$ le fibré principal des directions aux nœuds non ouverts et aux
points marqués, de groupe $\mathbb{U}^{B}$ où $B$ est ici $\hat A
\smallsetminus C$\footnote{La page pose $B = A \smallsetminus C$ à la page 105
mais écrit à la page 108 un isomorphisme de $\mathbb{U}^{\hat A}$-torseurs ;
les points marqués sont donc comptés dans $B$.}, de sorte que $SP_G = P_C \times_{M_G}
P_B$. Le changement de base par l'équivalence d'homotopie $M^C_G \to M_G$
donne $\widetilde{P}^B_G$ et $S\mathcal{T}_G = \pi_1(SP_G) \xleftarrow{\sim}
\pi_1(\widetilde{P}^B_G)$ ; et le carré

\begin{tikzcd}
\widetilde{P}^B_G \arrow[r] \arrow[d] & SP_{G'} \arrow[d] \\
\widetilde{M}^C_G \arrow[r] & M_{G'}
\end{tikzcd}

est cartésien : pour $X'$ construite à partir de $(X, (D_x))$, les directions
en $B$ ne changent pas. D'où le morphisme cherché
\[
S\mathcal{T}_G \simeq \pi_1(\widetilde{P}^B_G) \longrightarrow
\pi_1(SP_{G'}) = S\mathcal{T}_{G'} .
\]

Le monoïde multiplicatif $]0,1]^{C}$ agit sur $\widetilde{M}^C_G$ par
homothétie des disques, le même rayon aux deux branches d'un nœud, et le
morphisme $]0,1]^C \times \widetilde{M}^C_G \to M_{G'}$ joue le rôle de
l'inclusion dans $M_{G'}$ du voisinage tubulaire épointé de $M_G$ dans
$\widehat{M}_{G'}$\footnote{La page pose d'abord $\mu_{\vec C} =\,
]0,1]^{\vec C}$, puis ajoute « il vaut mieux prendre $]0,1]^{C}$ », ce que
confirme la marge ($\rho_{x'} = \rho_{x''}$).}. Si $M^C_G \to M_G$ a une
section (« y en a-t-il toujours ?? — je pense que oui… »), on a un carré
cartésien $SP_G \times\, ]0,1]^C \to SP_{G'}$ au-dessus de $P_C \times\,
]0,1]^C \to M_{G'}$, et $P_C \times\, ]0,1]^C$ s'identifie, « dans un sens
très fort », à une forme du voisinage tubulaire : le fibré en polydisques
épointés $\prod_{x \in C} (T_{x'} \otimes T_{x''})^{*}$.

Les cas les plus intéressants, conclut la page 111, sont $G' = G^\natural$, et
ceux où $G$ est de dimension modulaire $0$ ou $1$ et $G'$ de dimension
modulaire au plus $2$, notamment les généralisations immédiates
($\operatorname{card} C = 1$).

\pagerange{111}{116}
\subsection*{Dimension modulaire nulle : les points de base réels (pages 111 à
116)}

Si $G$ est de dimension modulaire nulle, $X$ est déterminée par $G$ à
isomorphisme unique près, et $M_G$ est un point. On choisit les disques
canoniquement. Sur $\mathbb{P}^1$ marquée par $0, 1, \infty$, de groupe
d'automorphismes $\mathfrak{S}_3$, on prend la métrique de la sphère ronde de
rayon $1$ invariante par $\mathfrak{S}_3$, les trois points formant un
triangle équilatéral sur un grand cercle $\Sigma_{\mathbb{R}}$, et autour de
chaque $i$ le cercle $\Gamma_{i\rho}$ de centre $i$ et de rayon rectiligne
$\rho$, assez petit pour que les trois disques soient disjoints\footnote{La
page dit $0 < \rho \leq 1$ puis demande le cercle « petit » ; la page 113
appelle $\Sigma_0$ le grand cercle nommé $\Sigma_{\mathbb{R}}$ à la page 112.
L'invariance du cercle par la transposition qui échange les deux autres
points est « automatique ».}. Le cercle $\Gamma_{i\rho}$ coupe
$\Sigma_{\mathbb{R}}$ en deux points diamétralement opposés, en bijection
canonique avec $I \smallsetminus \{i\}$ (le plus proche), d'où
\[
\Gamma_{i\rho} \simeq \mathbb{U} \wedge_{\{\pm1\}} (I \smallsetminus \{i\}).
\]
En chaque point marqué, on a donc deux directions réelles, celles qui
pointent vers les deux autres points.

Pour la courbe $X_G$, de normalisée $\coprod_s \Sigma_{\hat{\vec A}_s}$ recollée
selon $\sigma$, on a pour chaque arête $a$ un ensemble à deux éléments
$\varepsilon_a$ et $T_a \simeq \mathbb{U} \wedge_{\{\pm1\}} \varepsilon_a$ :
\[
\varepsilon_a = \hat{\vec A}_s \smallsetminus \{a\} \ \text{ si } a \in I, \qquad
\varepsilon_a = \bigwedge_{\tilde a \text{ sur } a} \bigl(\hat{\vec
A}_{\mathrm{or}(\tilde a)} \smallsetminus \{\tilde a\}\bigr) \ \text{ si } a \in A,
\]
le second étant canoniquement trivial quand $a$ est une boucle\footnote{Pour
une boucle en $s$, d'arcs $a', a''$, et de troisième arc $b$, les deux
ensembles $\{a'', b\}$ et $\{a', b\}$ ont une bijection canonique ($b \mapsto
b$), d'où la trivialisation $T_a \simeq \mathbb{U}$ de la page 113.}. D'où
\[
SP_G = \prod_{a \in \hat A} T_a = \mathbb{U}^{\hat A} \wedge_{\{\pm1\}^{\hat A}}
S\varepsilon_G, \qquad S\varepsilon_G = \prod_{a \in \hat A} \varepsilon_a ,
\]
et $S\varepsilon_G \subset SP_G$ est un ensemble de $2^{\operatorname{card}
\hat A}$ points de base\footnote{La page écrit « $\varepsilon_G$ (de cardinal
$2^{\operatorname{card} A}$) est contenu dans $SP_G$ », puis $S\varepsilon_G
\subset SP_G$ : le produit sur $\hat A$ a $2^{\operatorname{card} \hat A}$
éléments, celui sur $A$ (page 116) $2^{\operatorname{card} A}$.}. La
chirurgie, avec des rayons dans $]0,1[^{\vec A}$, donne un morphisme
$]0,1[^{\vec A} \times SP_G \to SP_{G^\natural}$, donc
$S\mathcal{T}_G = \pi_1(SP_G) \simeq \mathbb{Z}^{\hat A} \to
S\mathcal{T}_{G^\natural}$, et mieux un morphisme de groupoïdes fondamentaux ;
comme $]0,1[^{\vec A}$ est contractile, chaque point de $S\varepsilon_G$ donne
un point de base de $S\mathcal{T}_{G^\natural}$ — un point base « à l'infini »
de l'espace des courbes lisses\footnote{La page écrit $SPM_{G^\natural}$, les
lettres serrées ; c'est $SP_{G^\natural}$, comme sa parenthèse le dit.}. Soit
$S_0\mathcal{T}'_G$ le sous-groupoïde plein de $S\mathcal{T}_G$ induit sur
$S\varepsilon_G$ ; on a $S_0\mathcal{T}'_G \to S\mathcal{T}_{G^\natural}$.

\emph{Présentation.} Pour $\xi \in S\varepsilon_G$ et $\alpha \in
\{\pm1\}^{\hat A}$, soit $u_{\alpha, \xi} : \xi \to \alpha\xi$ le chemin qui
fait faire un demi-tour, $t \mapsto e^{i\pi t}\xi_a$, aux coordonnées où
$\alpha_a = -1$ et fixe les autres. Alors $u_{\beta, \alpha\xi}\, u_{\alpha,
\xi} = u_{\alpha\beta, \xi} + n(\alpha, \beta)$, où $n(\alpha,\beta)_a = 1$ si
$\alpha_a = \beta_a = -1$ et $0$ sinon (deux demi-tours font un tour, qui est
l'unité de $\mathbb{Z}^{\hat A}$). Sans utiliser la
$\mathbb{Z}^{\hat A}$-structure : générateurs $u_{a, \xi} : \xi \to \sigma_a
\xi$, pour $a \in \hat A$, où $\sigma_a$ change la seule coordonnée $a$, et
relations
\begin{enumerate}
\item[(1)] les carrés $u_{b, \sigma_a\xi}\, u_{a, \xi} = u_{a, \sigma_b\xi}\,
u_{b, \xi}$ pour $a \neq b$ ;
\item[(2)] les tours complets $v^a_\xi = u_{a, \sigma_a \xi}\, u_{a, \xi}$
commutent entre eux, $\xi$ fixé.
\end{enumerate}
Le groupoïde ainsi présenté est $S_0\mathcal{T}'_G$. C'est juste, et (2) est
même conséquence de (1) : les objets, les générateurs et les carrés (1) sont
les cellules de dimension $0$, $1$ et $2$ de la décomposition cellulaire
produit du tore $(\mathbb{U})^{\hat A}$ où chaque cercle a deux sommets
$\pm 1$, et le groupoïde fondamental d'un complexe cellulaire ne dépend que de
son $2$-squelette\footnote{Cette remarque et sa justification sont de cette
édition.}.

En prenant les seules arêtes liées, $\varepsilon_G = \prod_{a \in A}
\varepsilon_a \subset P_G$, on obtient de même ${}_0\mathcal{T}_G$, et le
carré commutatif
\[
S_0\mathcal{T}'_G \to S\mathcal{T}_{G^\natural}, \qquad {}_0\mathcal{T}_G \to
\mathcal{T}_{G^\natural}
\]
(flèches verticales d'oubli des points marqués). Le lien avec les
$\varepsilon(G)$ des pages 10 à 12 : une orientation de $\hat{\vec A}_s$
désigne, pour chaque arc en $s$, l'arc qui le suit, d'où une application
$\varepsilon(G) \to S\varepsilon_G$ compatible à $\{\pm1\}^S \to \{\pm1\}^{\hat
A}$, et un foncteur pleinement fidèle $S_0\mathcal{T}_G \to
S_0\mathcal{T}'_G$\footnote{Ce lien n'est pas écrit sur les pages ; il est de
cette édition. L'application $\varepsilon(G) \to S\varepsilon_G$ est injective
quand $G$ est connexe et a au moins un point marqué.}.

\pagerange{117}{124}
\subsection*{Dimension modulaire $1$ (pages 117 à 124)}

Tous les sommets sont de genre $0$ et d'ordre $3$, sauf un, $s_0$, qui est de
genre $0$ et d'ordre $4$, ou de genre $1$ et d'ordre $1$. La page traite le
premier cas. Soit $J = \hat{\vec A}_{s_0}$, à quatre éléments ; alors $M_G
\simeq M_{0,J} \simeq \mathbb{P}^1 \smallsetminus \{0, 1, \infty\}$ (après
numérotation), et $M_{0,J} = \mathrm{Rep}(J) \wedge_{\mathfrak{S}_4}
M^{!}_{0,4}$, où $\mathfrak{S}_4$ agit sur $M^{!}_{0,4} \simeq
\mathbb{P}^1 \smallsetminus \{0,1,\infty\}$ à travers $\mathfrak{S}_4 \to
\mathfrak{S}_3$, le groupe des permutations des trois partitions de type
$(2,2)$\footnote{$\mathrm{Rep}(J) = \mathrm{Bij}(\{1,2,3,4\}, J)$ ; la page
numérote les partitions par $P_i = \{\{i,4\}, \{1,2,3\} \smallsetminus
\{i\}\}$. Elle écrit $\mathbb{U}_{0,3}$ pour $\mathbb{P}^1 \smallsetminus
\{0,1,\infty\}$, et note $J$ d'une majuscule cursive à la page 117.}. Une
courbe de type $G$ équivaut à une droite projective $Y$ marquée par $J$.

Les arêtes qui ne touchent pas $s_0$, $\hat A_1$, ont leurs torseurs comme en
dimension nulle. Celles qui touchent $s_0$ sont de trois sortes, selon l'arc
$i \in J$ qui les porte : (I) une arête libre en $s_0$ ; (II) une boucle en
$s_0$, d'où une involution sans point fixe sur la partie $J_{\mathrm{II}}$ ;
(III) une arête de $s_0$ vers un autre sommet $s$, d'où, pour chaque $i \in
J_{\mathrm{III}}$, un ensemble à deux éléments $\tilde J_{\mathrm{III},i} = \hat{\vec
A}_s \smallsetminus \{\bar\imath\}$. Avec $J = J_{\mathrm{I}} \amalg
J_{\mathrm{II}} \amalg J_{\mathrm{III}}$,
\[
T^Y_{J_{\mathrm{I}}} = \prod_{i \in J_{\mathrm{I}}} T_i, \qquad
T^Y_{J_{\mathrm{II}}} = \prod_{\alpha \in J_{\mathrm{II}}/\sigma}
\ \bigwedge_{i \text{ sur } \alpha} T_i, \qquad
T^Y_{J_{\mathrm{III}}} = \prod_{i \in J_{\mathrm{III}}} T_i \wedge_{\{\pm1\}}
\tilde J_{\mathrm{III},i},
\]
où $T_i$ est le torseur tangent de $Y$ en $i$ ; ce sont des
$\mathbb{U}$-torseurs sur $M_{0,J}$, et
\[
SP_G \simeq T_{\hat A_1} \times T_{A_2}, \qquad A_2 = \hat A \smallsetminus
\hat A_1 \simeq J_{\mathrm{I}} \amalg J_{\mathrm{II}}/\sigma \amalg
J_{\mathrm{III}},
\]
$T_{\hat A_1}$ étant un $\mathbb{U}^{\hat A_1}$-torseur constant et
$T_{A_2}$ un $\mathbb{U}^{A_2}$-torseur sur $M_{0,J}$\footnote{La page passe de
$J_1, J_2$ (arcs liés non boucles, et les autres) à $J_{\mathrm{I}},
J_{\mathrm{II}}, J_{\mathrm{III}}$, et écrit tantôt $A_1$, tantôt $\hat A_1$ ; on
a unifié. Chaque $T_i$ figure une fois et une seule dans le produit, « sous
forme tordue » pour les arcs de type III.}.

\emph{Les points de base.} On les prend au-dessus des points de $M_{0,J}$ où
$(Y, J)$ a plus d'automorphismes que le groupe de Klein d'ordre $4$ qu'a toute
configuration de quatre points : les points de ramification de $M_{0,J} \to
M_{0,J}/\mathfrak{S}_3$, au nombre de cinq\footnote{« $5$ » est lu avec doute ;
c'est le compte juste, $3 + 2$. La page dit « des automorphismes d'ordre $> 4$
», c'est-à-dire un groupe d'automorphismes d'ordre $> 4$.}. Il y en a de deux
sortes.

a) Les trois configurations \emph{harmoniques}, où $J$ est l'ensemble des
sommets d'un carré, les automorphismes du carré étant induits par des
automorphismes projectifs (groupe d'ordre $8$). Une structure de carré $Q$ sur
$J$ équivaut à une partition de type $(2,2)$, celle des diagonales\footnote{La
page lit « partition de type $(1,2)$ », lecture incertaine ; c'est $(2,2)$.}.
Le groupe des rotations du carré agit simplement transitivement sur $J$ et
identifie les $T_i$ entre eux ; une orientation $\omega_Q$ du carré donne en
chaque sommet la direction vers le suivant, d'où $T_i \simeq \mathbb{U}
\wedge_{\{\pm1\}} \omega_Q$, et
\[
T^Y_{J_{\mathrm{I}}} = \mathbb{U}^{J_{\mathrm{I}}} \wedge_{\{\pm1\}} \omega_Q
\ (\text{diagonalement}), \qquad T^Y_{J_{\mathrm{II}}} \simeq
\mathbb{U}^{J_{\mathrm{II}}/\sigma}, \qquad T^Y_{J_{\mathrm{III}}} \simeq
\prod_{i \in J_{\mathrm{III}}} \mathbb{U} \wedge_{\{\pm1\}} (\omega_Q \wedge
\tilde J_{\mathrm{III},i}),
\]
le deuxième parce que le produit contracté de deux copies de $\mathbb{U}
\wedge \omega_Q$ est canoniquement $\mathbb{U}$\footnote{La page écrit
$T^Y_{J_{\mathrm{II}}} \simeq \mathbb{U}^{J_2}$ ; il y a un facteur par paire.}.
Si $J_{\mathrm{I}} \neq \emptyset$, on pose
\[
\varepsilon_{G,Q} = \Bigl(\prod_{a \in \hat A_1} \varepsilon_a\Bigr) \times
\omega_Q \times \prod_{i \in J_{\mathrm{III}}} (\omega_Q \wedge \tilde
J_{\mathrm{III},i}),
\]
torseur sous $\gamma_G = \{\pm1\}^{\hat A_1} \times \{\pm1\} \times
\{\pm1\}^{J_{\mathrm{III}}}$, qui s'envoie dans $\mathbb{U}^{\hat A}$ par les
inclusions évidentes et, sur le facteur du milieu, diagonalement dans
$\mathbb{U}^{J_{\mathrm{I}}}$ ; et l'on a $SP^Y_G \simeq \mathbb{U}^{\hat A}
\wedge_{\gamma_G} \varepsilon_{G,Q}$ ; si $J_{\mathrm{I}} = \emptyset$, on omet
le facteur $\omega_Q$\footnote{La page note ces ensembles $\mathcal{E}_{G,Q}$ ;
on garde $\varepsilon$ pour les ensembles de points de base.}. Le groupoïde
$S_0\mathcal{T}^Q_G$ est le sous-groupoïde plein du groupoïde fondamental basé
en $\varepsilon_{G,Q}$ ; ses formules explicites par générateurs et relations
sont renvoyées à plus tard (« j'y reviendrai »).

b) Les deux configurations \emph{équianharmoniques} (page 124), où les
automorphismes de $(Y, J)$ fixant un point $i_0$ forment un groupe d'ordre $3$
(le groupe total, $\mathfrak{A}_4$, est d'ordre $12$). Les permutations
circulaires des trois autres points forment un $\mu_3$-torseur $\eta_{i_0}$, et
le torseur tangent en $i_0$ est $\mathbb{U} \times_{\mu_3} \eta_{i_0}$, par
l'inclusion $\mu_3 \hookrightarrow \mathbb{U}$\footnote{La page écrit « $J_i
\simeq \mathbb{U} \times_{\mu_3} \eta_i$ ($i \in I$) » ; il s'agit du torseur
tangent en $i$. Le paragraphe qui précède est très raturé.}. La page s'arrête
sur « On ».

\pagerange{125}{132}
\section*{Graphes MD (pages 125 à 132)}

La chemise de la page 125 ouvre six pages qui fixent la combinatoire utilisée
depuis la page 2.

\emph{Le tableau (page 126).} Pour chaque type $(g, \nu)$ de petite
dimension, les graphes MD du type, reliés par des flèches de spécialisation
vers les plus dégénérés, certaines affectées d'une multiplicité ($2$, $3$), et
en regard la dimension $3g - 3 + \nu$ : $0$ pour $(0,3)$, $1$ pour $(0,4)$ et
$(1,1)$, $2$ pour $(0,5)$ et $(1,2)$, $3$ pour $(0,6)$, $(1,3)$ et $(2,0)$, $4$
pour $(2,1)$, ce dernier inachevé.

\emph{Les définitions (pages 127 et 128).} Un graphe pondéré est un quintuple
$(S, \hat{\vec A}, \mathrm{or}, \sigma, g)$ comme dans les conventions ci-dessus ;
la condition de stabilité est
\[
(4)\qquad 2g_s + \hat\nu_s \geq 3 \quad \text{pour tout } s \in S,
\]
soit $\hat\nu_s \geq 3$ si $g_s = 0$, $\hat\nu_s \geq 1$ si $g_s = 1$, rien
si $g_s \geq 2$. On pose $\nu(G) = \operatorname{card} I$, l'ordre marqué ;
$c(G) = \operatorname{card} A$, le nombre d'arêtes liées, \emph{codimension
modulaire} ; et le \emph{genre virtuel}
\[
(8)\qquad g = \sum_{s \in S} g_s + h^1(G),
\]
qui est $\dim H^1(X, \mathcal{O}_X)$ pour une courbe de type $G$, $G$ connexe.
La dimension du type est $3g - 3 + \nu$, et celle de la strate est
\[
\delta(G) = \dim M_G = \sum_{s} (3g_s - 3 + \hat\nu_s) = 3g - 3 + \nu - c(G),
\]
d'où
\[
(10)\qquad \delta(G) + c(G) = 3g - 3 + \nu .
\]
La formule (9) de la page est donc la dimension du type\footnote{La page définit (9)
$\delta(G) = 3g - 3 + \nu$, « dimension modulaire », puis écrit (10) $\delta(G)
+ c(G) = \delta(\mathbf{G})$, le second membre d'une autre lettre, le type. La
formule (9) est donc celle du type, et la dimension modulaire du graphe, celle
que (10) et les pages 111 et 117 utilisent, est $3g - 3 + \nu - c(G)$ : le
calcul de la somme utilise $h^1(G) = c(G) - \operatorname{card} S + 1$. Les
équations (5) à (7) confondent $\nu_s$ et $\vec\nu_s$ ; (6) donne $\nu =
\operatorname{card}(\vec A^{\bar\sigma})$, c'est $\operatorname{card} I$.}.

\emph{La contraction.} Un morphisme $u : G \to G'$ (que la page appelle
homomorphisme de spécialisation, et les pages 103 et 105 de généralisation :
une courbe de type $G$ est une dégénérescence des courbes de type $G'$) est un
couple $u_s : S \to S'$, $u_a : \hat{\vec A} \to \hat{\vec A}' \amalg S'$ tel
que $\vec B = u_a^{-1}(\hat{\vec A}')$ contienne $I$ et soit stable par
$\sigma$, que $u_a|_{\vec B}$ commute aux involutions et aux origines, que les
arcs contractés $\vec C = u_a^{-1}(S')$ aillent sur le sommet image de leur
origine, que $S/R \to S'$ et $\vec B \to \hat{\vec A}'$ soient bijectifs ($R$
la relation d'équivalence engendrée par les arêtes contractées), et que le
genre de $s'$ soit le genre virtuel du sous-graphe contracté sur $s'$. Le
graphe $G'$ se déduit alors de $G$ et de $C = \vec C/\sigma$ à isomorphisme
unique près :
\[
S' = \pi_0(G_{0C}), \quad \hat{\vec A}' = \hat{\vec A} \smallsetminus \vec C,
\quad \mathrm{or}' = \text{la composée } \hat{\vec A}' \to S \to S', \quad
\sigma' = \sigma|, \quad g_{s'} = g(G_{0C,s'}),
\]
où $G_{0C}$ est le sous-graphe pondéré de sommets $S$ et d'arcs $\vec C$ ; on
a $I \simeq I'$ et $\pi_0(G) \simeq \pi_0(G')$, avec les mêmes genres
virtuels par composante. C'est la contraction des graphes stables
d'aujourd'hui, celle qui décrit l'adhérence des strates\footnote{La page
répète « est une bijection » dans d) ; la condition d) équivaut, selon la
marge, à la connexité des $G(s')$.}.

\emph{Les opérations (pages 130 et 132).} 1) Le « décampage » d'arêtes et le
recollement d'arêtes libres correspondent à la normalisation partielle et au
recollement ; $S$, les genres et les ordres totaux sont conservés, et $M_G
\simeq M_{G'}$ : c'est ce qui sert à écrire $M_G$ comme le produit des
$M_{g_s, \hat{\vec A}_s}$ (pages 2 et 10). 2) Les homomorphismes ordinaires de
graphes pondérés, compatibles aux genres, et les morphismes $M_G \to M_{G'}$
qu'ils induisent — ceux de transition entre les $M_{g,\nu}$\footnote{La page
barre d'une croix les variantes pour les compactifiés, avec « ne fait pas
sens ».}. 3) Les homomorphismes d'\emph{induction} $u : G' \to G$, compatibles
aux origines et aux genres, injectifs sur les arcs de chaque sommet, envoyant
$\vec A'$ dans $\vec A$ en commutant aux involutions, mais sans
compatibilité aux extrémités : d'une courbe $X$ de type $G$, de normalisée
$\coprod_s \widetilde{X}_s$, on tire une courbe $X'$ de type $G'$ en posant
$\widetilde{X}'_{s'} = \widetilde{X}_{u(s')}$, marquée par les arcs de $s'$, et en
recollant selon $\sigma'$ ; on a un morphisme fini et net (non ramifié) $X'
\to X$, qui envoie $I'$ dans $I \amalg \mathrm{Sing}\,X$. Le dossier s'arrête
sur la réciproque : « Inversement \ldots{} un tel morphisme ».

\end{document}
